%% The Visual Computer submission draft.
%% Generic Springer-compatible LaTeX layout with 10 pt Times, A4 paper,
%% decimal section numbering and continuous page numbers.
%% Source of record for the numbers: _deposit/ssgf-llvip-fusion-detection-v1.
%
\documentclass[10pt,a4paper]{article}
\usepackage[T1]{fontenc}
\usepackage{mathptmx}
\usepackage[a4paper,top=2.5cm,bottom=2.5cm,left=2.5cm,right=2.5cm]{geometry}
\usepackage{amsmath,amssymb,graphicx,booktabs,siunitx}
\usepackage{etoolbox}
\usepackage[hidelinks]{hyperref}

\AtBeginEnvironment{center}{\setlength{\topsep}{2pt}%
  \setlength{\partopsep}{0pt}}

% Tables: one font size everywhere. Tabulars are set at \small and the column
% separation is fixed, so a wide table is fitted by its own column widths rather than
% by scaling the whole table, which would give every table a different apparent size.
\setlength{\tabcolsep}{4pt}
\AtBeginEnvironment{tabular}{\small}
\AtBeginEnvironment{tabular*}{\small}
% Tighten the glue around floats and the reference list. With eleven figures and
% fourteen tables this keeps the spacing uniform between displays.
% spacing uniform between displays.
\setlength{\floatsep}{7pt plus 2pt minus 2pt}
\setlength{\textfloatsep}{9pt plus 2pt minus 3pt}
\setlength{\intextsep}{7pt plus 2pt minus 2pt}
\setlength{\dblfloatsep}{4pt plus 1pt minus 1pt}
\setlength{\dbltextfloatsep}{5pt plus 1pt minus 1pt}
\setlength{\abovecaptionskip}{2pt}
\setlength{\belowcaptionskip}{0pt}
\AtBeginEnvironment{thebibliography}{\setlength{\itemsep}{0pt}%
  \setlength{\parsep}{0pt}\setlength{\topsep}{2pt}}

\hyphenation{SPGFu-sion PFusion}

\begin{document}

\title{Fidelity Metrics Do Not Reliably Rank Visible--Infrared Fusion Rules
for Pedestrian Detection}

\author{Jianyong Zhao \and Zigao Wu \and Mengfei Ma\\[4pt]
\small North China Electric Power University, Baoding 071003, China}
\date{}

\maketitle
\begingroup
\renewcommand{\thefootnote}{}
\footnotetext{Corresponding author: Zigao Wu (zgwu@ncepu.edu.cn). ORCID:
J. Zhao 0009-0005-3251-2785; Z. Wu 0000-0002-7449-343X;
M. Ma 0009-0007-1328-1314.}
\endgroup

\begin{abstract}
Fidelity metrics are widely used to rank visible-infrared fusion rules before a
downstream detector is trained, but the ranking can depend on the evaluator and on how
the metric is implemented. We keep fourteen fused inputs fixed and compare their
downstream pedestrian-detection performance under three independently tested
visible-trained detectors and a modality-agnostic detector. With all 300 LLVIP test
images held out, all 42 visible-trained evaluator-input contrasts are negative at
mAP50, with deficits from $0.046$ to $0.246$. The modality-agnostic evaluator reverses
the sign for four inputs by $0.022$ to $0.029$ in every run. Two constructions of the
same VIF family rank SPGFusion first and last, respectively. Paired bootstrap
intervals, four independent training runs, both source conventions and an
alignment-stratified analysis are used to distinguish stable effects from evaluator,
index, sampling and training variability. The results show that fidelity scores should
not be used as a sole proxy for detector utility. The conclusion is bounded to the
measured pedestrian-detection populations, and source alignment remains a material
source of variability.
\end{abstract}

\noindent\textbf{Keywords:} visible-infrared image fusion; image-quality assessment;
pedestrian detection; evaluation metrics; benchmark reliability; multimodal perception

\section{Introduction}

Visible-infrared image fusion combines complementary sensing modalities into an
intermediate image used by downstream perception systems. In pedestrian detection, the
infrared source preserves thermal targets while the visible source supplies texture and
appearance. Different fusion rules redistribute trust between the two sources, and the
selected rule can materially change what a detector sees. Because retraining a detector
for every candidate rule is expensive, fusion rules are commonly ranked by
full-reference fidelity metrics before a task-specific evaluation is performed.

Infrared and visible images are fused so that a thermal target survives into one
representation while visible texture is retained. The rules differ in which modality
they trust at each pixel. Because the downstream task is expensive to retrain, the rule
is usually selected by fidelity metrics such as VIF, $Q^{\mathrm{AB/F}}$ or SSIM.
VIF compares source information retained in the fused image with information lost to a
modelled noise channel across spatial scales. $Q^{\mathrm{AB/F}}$ scores gradient
transfer from the two sources, and SSIM combines luminance, contrast and structure. All
three are higher-is-better. Their agreement with human judgement and with one another
is incomplete~\cite{ma2019}. Recent reviews of infrared--visible fusion organize the
field around visual quality, benchmarks and evaluation protocols~\cite{gajjar2026visual},
but the detector-dependent reliability of a fidelity ranking remains difficult to
infer from the fused image alone.

We therefore hold the fused images fixed and ask whether the ordering induced by a
fidelity score agrees with downstream detection performance. The same inputs are
evaluated under independently trained detectors, two source conventions and two
constructions of the same VIF family. Repeated training, paired image resampling and an
alignment-stratified analysis are used to distinguish robust effects from evaluator,
index, sampling and training variability. Figure~\ref{fig:framework} summarizes the
fusion and downstream evaluation pipeline.

\begin{figure*}[!t]\centering

\includegraphics[width=\linewidth]{framework_v7.pdf}
\caption{Overall framework. The shared encoder $E$ maps both sources to feature stacks
$F_v$ and $F_r$. A per-pixel targetness map $S_s$, either the morphological top-hat
prior or the learned probe $\Phi$, is multiplied by the feature activity $A_s$ into
$d_s=S_s\odot A_s$; a temperature softmax with $\tau=0.3$ turns the two maps into the
per-source weights $w_s$ of Eq.~\eqref{eq:rule}, and the mirrored decoder $D$
reconstructs the fused image. The evaluation that this article is about sits downstream
of $I_f$: its ranking depends on the evaluator, not only on the rule. The frame and the
two inset maps are illustrative outputs for LLVIP test image 190003; the maps show the
learned-probe targetness $S_v$ and the normalized visible-source weight $w_v$, and the
thin orange boxes mark the pedestrian annotations.}
\label{fig:framework}

\end{figure*}

We test fourteen fused variants on two pedestrian-detection benchmarks against 26
detection checkpoints, and repeat the comparison with one segmentation checkpoint. The
study makes three contributions.

\begin{itemize}\setlength\itemsep{0pt}
\item[(i)] \textbf{A detector-dependent evaluation.} We keep the fused inputs fixed and
compare fidelity rankings with downstream detection across independently trained
detectors, source conventions and metric constructions.
\item[(ii)] \textbf{Evidence that the ranking is evaluator- and construction-dependent.}
On the clean split, every fused input loses to the visible source under all three
visible-trained evaluators, with deficits from $0.046$ to $0.246$ mAP50. A
modality-agnostic reference reverses the sign for four inputs in every run. Two
constructions of the same VIF family place SPGFusion first and last. These are changes
in the conclusion drawn from the same fused inputs, not small changes in the magnitude
of a fixed ranking.
\item[(iii)] \textbf{A reproducible evaluation protocol.} Paired bootstrap intervals,
four independent training runs, two source conventions and an alignment-stratified
sensitivity analysis are released with the per-image scores, fused images and analysis
scripts.
\end{itemize}

Fusion rules range from hand-crafted multi-scale transforms, sparse representations and
edge-preserving decompositions, applied to infrared--visible, medical and multi-focus
pairs~\cite{zhu2017,zhang2021,yin2013,nejati2015,yin2015,wang2014,zhang2016,wei2017,fu2009,ma2021,chen2007,piella}.
Adversarial and generative designs~\cite{ma20191,hou2021,yi2021,song2022} and
unsupervised encoder--decoder networks~\cite{xu2022,li2021,li2019,zhang2020,wang2024,liu2023,hao2022,yang2023,an2024}
followed them, and the most recent rules are diffusion and transformer
models~\cite{yi2024,yang2025}. A subset of these
rules is task-aware: PIAFusion conditions on illumination~\cite{tang2022}, SPGFusion
injects a semantic prior~\cite{zhang2026}, MTFusion is guided by a segmentation network
~\cite{zhou2026}, and saliency-guided rules weight the sources by an importance
map~\cite{bianco2020,khan2025,wang2025,xu2021}. Fidelity is then read off reference and
non-reference indices~\cite{sheikh2006,ma2019}, and the saliency literature that
supplies the importance maps covers generic, contrast-based and foundation-model
saliency~\cite{qin2020,zhu2025,kirillov2023sam,bianco2020,guo2024,petilli2022,li2010}.

The present study asks whether a fidelity ranking survives a change in the downstream
perception task's evaluator and in the implementation of the metric. We evaluate this
ordering with detectors trained for thermal pedestrian detection
~\cite{redmon2016yolo,jocher2024yolo,lv2024rtdetr} and a segmentation network
~\cite{ronneberger2015} on public benchmark releases~\cite{flir_adas}.

Section~\ref{sec:validity} defines the evaluation design.
Section~\ref{sec:method} describes the fusion rules and experimental protocol.
Section~\ref{sec:fidelity} reports fidelity measurements, and
Section~\ref{sec:detect} compares them with detector performance.
Section~\ref{sec:ranking} separates evaluator and index effects.
Section~\ref{sec:adapt} tests task-aware adaptation.
Section~\ref{sec:discussion} states the practical limits of the result.

\section{Evaluation Design}
\label{sec:validity}

For a fixed set of fused inputs $\mathcal{I}$ and downstream evaluator $e$, let
$Y_e(i)$ denote the detection performance of input $i$ on a test set excluded from
training and checkpoint selection. Let $S_c(i)$ denote a fidelity score computed under
metric construction $c$. A user who selects $\arg\max_i S_c(i)$ assumes that the
ordering induced by $S_c$ agrees with that induced by $Y_e$.

We test this assumption in two ways. First, we compare the sign of each fused input's
contrast against the visible source under $S_c$ and $Y_e$. Second, we compare the full
rank order of the fourteen fused inputs. This analysis is repeated under different
evaluators, source conventions and VIF constructions.

Three sources of variability are reported separately. Sampling variability is estimated
by a paired image bootstrap that applies the same resampling index to both arms of a
contrast. Training variability is represented by the full range over four independently
trained YOLO11n runs. Evaluator and index variability are represented by changes in
sign or rank when the detector training modality or VIF construction changes. Source
alignment is examined through quartile and sequence sensitivity analyses because the
available ECC residual is translational and does not provide ground-truth
correspondence.

\section{Fusion Rules and Experimental Protocol}
\label{sec:method}

\begin{sloppypar}
Given a registered pair $(I_v,I_r)$, fusion produces $I_f$ that retains the visible
texture and the infrared thermal targets. The decision is which modality to keep at
each pixel.

A shared stride-1 encoder $E$ (reflection padding, Leaky-ReLU, $5\times5$ kernels,
channel widths $1,16,24,24,24$) maps each source to a 24-channel feature stack
$F_s$, and a mirrored decoder $D$ (widths $24,16,16,16,1$, final $\tanh$) reconstructs
an image from a weighted feature stack. Both are trained once, self-supervised, to
reconstruct a single source image,
\begin{equation}
\mathcal{L}_{\mathrm{rec}} = \lambda\,\mathrm{MSE}(x,\hat{x}) + (1-\mathrm{SSIM}(x,\hat{x})),\qquad \lambda=25,
\end{equation}
on $128\times128$ patches of the MSRS training set, so that the features, not the
fusion rule, carry the reconstruction quality.

Fusion reweights those features. Writing $A_s=\sum_c|F_s|$ for the feature activity and
$S_s\in[0,1]^{H\times W}$ for a per-pixel targetness map of source $s$, the
un-normalised weight is $d_s=S_s\odot A_s$, normalised by the shared per-image maximum
$m=\max(\max_x d_v,\max_x d_r)$ and passed through a temperature softmax,
\begin{equation}
\begin{aligned}
w_s&=\frac{\exp\!\big(d_s/(m\tau)\big)}
{\exp\!\big(d_v/(m\tau)\big)+\exp\!\big(d_r/(m\tau)\big)},\\
\hat{F}&=\sum_s w_s\odot F_s,\qquad I_f=D(\hat{F}),
\end{aligned}
\label{eq:rule}
\end{equation}
with $\tau=0.3$ throughout. The product $m\tau$, not $\tau$, sets how soft the rule
is, which is what makes fidelity comparisons between rules with different weight
spreads meaningful.

Two targetness maps instantiate $S_s$. The \emph{top-hat prior} is the morphological
white/black top-hat response, max-combined across the two polarities, smoothed by a
Gaussian of $\sigma=3$ and min--max normalised per image; it is parameter-light and
target-oriented. The \emph{learned probe} replaces it with a two-layer convolutional
head on the frozen encoder features, trained by distillation against the prior,
\begin{equation}
\mathcal{L}_{\mathrm{probe}}=\mathrm{BCE}\big(\sigma(\Phi(F)),\,S^{\mathrm{th}}\big),
\end{equation}
so that the same targetness idea becomes data-driven. The probe costs 0.07\,M
parameters and 19.1\,ms per $512\times384$ pair, orders of magnitude less than the
largest end-to-end network in the comparison. An image-level \emph{ensemble} of the
two rules is also reported. Figure~\ref{fig:framework} summarises the pipeline and
Supplementary Sec.~20 gives the layer configuration and the algorithm listing.
\end{sloppypar}

\begin{sloppypar}
The variant set isolates the parts of the rule. \emph{mean fusion} averages the
sources; the \emph{activity-only} rule sets $S_v=S_r\equiv1$, which reduces
Eq.~\eqref{eq:rule} to the classical $L_1$-norm activity weighting and is the control
that separates targetness from activity; \emph{spatially permuted} variants shuffle
each targetness map, preserving its marginal distribution while destroying its
alignment with the image; and \emph{weight-shuffled} variants apply one common spatial
permutation to $w_v$ and $w_r$, preserving both marginals and the complementarity
$w_v+w_r=1$ while destroying the alignment. IFCNN~\cite{zhang2020} and
DenseFuse~\cite{li2019} are retrained on the same data with their published
$L_1$-norm fusion, while PIAFusion~\cite{tang2022} and SPGFusion~\cite{zhang2026} are
reproduced with their official weights and protocols.
\end{sloppypar}

\paragraph{Evaluation protocol.} Detection uses 26 independently trained checkpoints
(Table~\ref{tab:inventory}): three visible-trained families on LLVIP under the pooled
sensitivity protocol (4 YOLO11n, 4 YOLO26n and 3 RT-DETR-l runs), one pooled
modality-agnostic run, the clean-split versions of those four evaluators
(4/1/1/4 runs), one larger YOLO11s backbone, and three FLIR detectors. A separate U-Net
is used for segmentation. The clean split and the pooled split are reported separately
throughout.
Each detector is applied to every fused input and to the two unfused sources at a
confidence threshold of $0.001$; detections are matched to
the ground truth by a one-to-one maximum-cardinality assignment at IoU 0.5, after sorting
each image's boxes by confidence. The true-positive flags are pooled over images, and the
average precision is the 101-point interpolation of the detector library (Ultralytics,
release 8.4.189). mAP50 is AP at IoU 0.50, mAP50-95 averages AP over IoU thresholds
$0.50{:}0.05{:}0.95$, and mIoU is the class-averaged intersection-over-union. In the
fidelity tables EN is image entropy, PSNR is peak signal-to-noise ratio, and SSIM is
reported separately against the visible and infrared sources. A contrast is called
\emph{resolved} when its 95 percent interval excludes zero. Every interval in this article resamples the same 300 images
1000 times and applies one draw to both arms of a contrast, so a reported difference
is paired rather than a difference of two independent scores. For each single-run
clean visible-only evaluator, Benjamini--Hochberg correction is applied separately to
the fourteen fused-input contrasts against its own visible source. The
\emph{reported protocol} scores every input, including the visible source, in a single
channel; the \emph{colour-corrected protocol} and its audit are in
Supplementary Sec.~14, and Sec.~\ref{sec:detect} states each claim in the convention
that supports it. Fidelity uses the index of Sec.~\ref{sec:fidelity} together with the
reference pixel VIF construction for the comparison of Sec.~\ref{sec:ranking}.

LLVIP contains 15,488 registered pairs in official partitions of 12,014 training and
3,474 test pairs. We retain the official partition and take 1,500 training and 300 test
pairs that contain at least one pedestrian box. The 300 test images contain 744 annotated
pedestrians. Their per-image counts are 98 images with one, 90 with two, 54 with three,
25 with four, 15 with five, 8 with six, 2 with seven, 5 with eight and 3 with nine. The
registered release is used as distributed and no additional spatial transform is applied;
LLVIP's raw unregistered pairs form a separate release. As an alignment sensitivity check,
we estimate each test pair's residual translation with ECC on normalized Sobel-edge maps
(median 2.53 px; best-quartile cutoff 1.61 px). This estimate is not used to generate or
score any fused image.

All primary LLVIP evaluators use the same fixed, test-independent partition. From the
official training partition we take the first 1,500 non-empty annotations in stem order,
assign the first 1,350 to training and the remaining 150 to validation, and exclude all
300 scored images from both. The three intersections among train, validation and test
contain zero stems; \texttt{split\_manifest.json} records every stem. The clean
YOLO11n and YOLO26n visible-only evaluators were started from their official checkpoints
and trained for 40 epochs at $640$ px with batch size 16 and seed 0;
YOLO11n was repeated for seeds 1 to 3. The clean RT-DETR-l evaluator used the same
images, epochs and seed, with batch size 4 to fit the available GPU. We selected the
\texttt{best.pt} checkpoint by the detector library's validation fitness on the 150
held-out validation pairs. The modality-agnostic evaluator is a YOLO11n trained on both
registered modalities of the same 1,350 pairs with the same recipe and four independent
seeds, concatenating the visible and infrared sources as channels and applying no
post-training fusion rule. The test images are not used for optimisation, checkpoint
selection, threshold selection or calibration.
The aligned FLIR release contributes 4,129 training and 1,013 validation pairs, from
which the same selection gives 1,500 and 300.

\begin{table*}[!t]\centering\small
\caption{Evaluator inventory. ``Pooled'' uses 1,500 training images and overlaps the
300 scored test images; ``clean'' uses the fixed 1,350/150/300 split. Runs count
independently trained checkpoints. The table contains 26 detection checkpoints plus
one segmentation checkpoint.}\label{tab:inventory}
\setlength{\tabcolsep}{3pt}
\begin{tabular*}{\linewidth}{@{\extracolsep{\fill}}lll c l l@{}}
\toprule
Benchmark & Evaluator & Input & Runs & Protocol & Role\\
\midrule
LLVIP & YOLO11n & visible & 4 & pooled & sensitivity\\
      & YOLO26n & visible & 4 & pooled & sensitivity\\
      & RT-DETR-l & visible & 3 & pooled & sensitivity\\
      & YOLO11n & VIS+IR & 1 & pooled & reversal sensitivity\\
\addlinespace
      & YOLO11n & visible & 4 & clean & primary\\
      & YOLO26n & visible & 1 & clean & primary\\
      & RT-DETR-l & visible & 1 & clean & primary\\
      & YOLO11n & VIS+IR & 4 & clean & primary reversal\\
\addlinespace
      & YOLO11s & visible & 1 & pooled & backbone sensitivity\\
FLIR  & YOLO11n & visible & 1 & independent & replication\\
      & YOLO26n & visible & 1 & independent & replication\\
      & YOLO11n & VIS+IR & 1 & independent & replication\\
MSRS  & U-Net & all inputs & 1 & independent & segmentation\\
\bottomrule
\end{tabular*}
\end{table*}

\section{Fidelity Measurements}
\label{sec:fidelity}

Table~\ref{tab:sota} compares fidelity against published methods and
Table~\ref{tab:fusion} against our four-rule family on three benchmarks.

\begin{table*}[!t]\centering\small

\caption{Fusion-quality VIF compared with published baselines (higher is
better), 512\,px, means over the test pairs. For each dataset, the first row is the
reported-index VIF and the second is the reference pixel VIF. Best in each row is
bold.}\label{tab:sota}
\setlength{\tabcolsep}{3pt}
\begin{tabular}{@{}ll SSSSSSS@{}}
\toprule
Dataset & Index & {IFCNN} & {DenseFuse} & {PIAFusion} & {SPGFusion} & {Ours (probe)} & {Ours (top-hat)} & {Mean}\\
\midrule
TNO & reported & 0.4440 & 0.3697 & 0.5485 & \textbf{0.7855} & 0.3297 & 0.3560 & 0.2683\\
    & reference & 0.2896 & 0.2729 & \textbf{0.3256} & 0.1607 & 0.2309 & 0.2307 & 0.2226\\
\addlinespace
MSRS & reported & 0.3757 & 0.3004 & 0.7507 & \textbf{0.8314} & 0.4015 & 0.3914 & 0.2914\\
     & reference & 0.2517 & 0.2722 & \textbf{0.4468} & 0.2308 & 0.2562 & 0.2573 & 0.2558\\
\addlinespace
RoadScene & reported & 0.4988 & 0.3828 & 0.6315 & \textbf{0.8293} & 0.4229 & 0.4240 & 0.3126\\
          & reference & \textbf{0.3082} & 0.3052 & 0.2855 & 0.2166 & 0.2602 & 0.2617 & 0.2661\\
\bottomrule
\end{tabular}

\end{table*}

\begin{table*}[!t]\centering\small

\caption{Full-benchmark fusion metrics (TNO=74, MSRS=361, RoadScene=221,
512\,px). Values are means over the test images; the per-image standard deviation is
roughly 0.02 to 0.08 in VIF across methods. The VIF values are the same measurements as
the corresponding entries in Table~\ref{tab:sota}; the remaining columns report
additional metrics. Best in bold.}\label{tab:fusion}
\setlength{\tabcolsep}{4pt}
\begin{tabular*}{\linewidth}{@{\extracolsep{\fill}}l l SSSS S@{}}
\toprule
Dataset & Method & {Rep. VIF} & {EN} & {SSIM\textsubscript{vis}} & {SSIM\textsubscript{ir}} & {PSNR\textsubscript{vis}}\\
\midrule
TNO & mean & 0.2683 & 5.9972 & 0.8058 & 0.7433 & 21.2620\\
    & \textbf{top-hat} & \textbf{0.3560} & \textbf{6.1219} & 0.8193 & 0.7172 & \textbf{21.7204}\\
    & probe & 0.3297 & 6.0719 & 0.8115 & 0.7323 & 21.3373\\
    & ensemble & 0.3431 & 6.0960 & 0.8167 & 0.7251 & 21.5370\\
\addlinespace
MSRS & mean & 0.2914 & 5.7475 & 0.7418 & 0.6538 & 18.1703\\
    & top-hat & 0.3914 & 5.7813 & 0.7384 & 0.6582 & 18.2683\\
    & \textbf{probe} & \textbf{0.4015} & \textbf{5.8080} & \textbf{0.7545} & 0.6452 & \textbf{18.7262}\\
    & ensemble & 0.3955 & 5.7938 & 0.7473 & 0.6518 & 18.5072\\
\addlinespace
RoadScene & mean & 0.3126 & 7.1115 & 0.6868 & 0.7216 & 16.4718\\
    & \textbf{top-hat} & \textbf{0.4240} & \textbf{7.2220} & 0.6653 & 0.7283 & 15.8169\\
    & probe & 0.4229 & 7.1974 & 0.6685 & 0.7258 & \textbf{16.2372}\\
    & ensemble & 0.4234 & 7.2055 & 0.6680 & 0.7283 & 16.0380\\
\bottomrule
\end{tabular*}

\end{table*}

\begin{sloppypar}
SPGFusion and PIAFusion dominate pixel fidelity, while IFCNN exceeds DenseFuse and both
targetness rules on TNO but not either end-to-end network. This is the expected
outcome for large energy-trained models, and the learned probe does not exceed them on
any benchmark. Within our own family the target-aware top-hat prior beats mean fusion in VIF
on every dataset ($+33\%$, $+34\%$ and $+36\%$), and the probe is the best of the
family on MSRS while trailing the prior slightly on TNO and RoadScene. The latter is not
an independent failure of the learned rule: the probe is distilled from the prior, so
near-parity with its teacher is the expected outcome and a blanket gain over the teacher
is not required by the design. The MSRS reversal shows that learned features can move
the fidelity ranking, but that movement does not translate into a downstream gain. The
classification-saliency rule of Xu et al.~\cite{xu2021}, run through its own
configuration, sits within $0.018$ of mean fusion on all three benchmarks
($+0.0101$ on TNO, $-0.0043$ on MSRS, $-0.0177$ on RoadScene, with paired intervals
excluding zero in each case), so a modality-cue saliency buys about one hundredth of a
VIF point over mean fusion. That is the fidelity story, and the rest of the article is
about what it is worth downstream. Supplementary Tables~S1 to S3 give the remaining
metrics and the component ablation, and Supplementary Secs.~1 to 6
give the sharpness control, the saliency diagnostics and the qualitative comparisons.

Two constructions of the same named index are used. The \emph{reported} index
accumulates four-scale VIF with Gaussian widths $2^{s-1}$, no decimation between scales,
and the numeric noise variance $2.0$ applied to inputs normalized to $[0,1]$. The
\emph{reference} pixel VIF uses the canonical widths $n/5$, decimation between scales,
and the noise variance $2/255^2$ in those same input units. Every fidelity number in
this article is the reported construction unless it is labelled otherwise;
Sec.~\ref{sec:ranking} shows that the two place SPGFusion first and last on the same
images, including on all three fidelity benchmarks. Table~\ref{tab:cost} measures all
four author rules and the four reproduced baselines on the same device and input size.

\begin{table*}[!t]\centering\footnotesize

\caption{Inference cost at $512\times384$, FP32, batch size one, on an NVIDIA
RTX 2080 Ti. Latency is averaged over repeated forward passes after three warm-up passes;
the SPGFusion entry uses five passes after two warm-ups because of its cost. Extra peak
memory is allocated GPU memory above the pre-call baseline. SPGFusion parameters count
the fusion network only; its memory also includes the DINO and CLIP feature extractors
required at inference.}\label{tab:cost}
\setlength{\tabcolsep}{4pt}
\begin{tabular*}{\linewidth}{@{\extracolsep{\fill}}l SSS@{}}
\toprule
Method & {Params (M)} & {Latency (ms)} & {Extra peak memory (MiB)}\\
\midrule
mean fusion & 0.06 & 12.4 & 120.3\\
top-hat prior & 0.06 & 21.2 & 132.1\\
learned probe & 0.07 & 19.1 & 149.3\\
ensemble & 0.07 & 41.0 & 188.3\\
IFCNN~\cite{zhang2020} & 0.08 & 7.0 & 242.2\\
DenseFuse~\cite{li2019} & 0.04 & 6.3 & 182.2\\
PIAFusion~\cite{tang2022} & 1.18 & 39.2 & 869.1\\
SPGFusion~\cite{zhang2026} & 75.86 & 538.5 & 1454.8\\
\bottomrule
\end{tabular*}

\end{table*}
\end{sloppypar}

\section{Task-Aware Validity of Fidelity Rankings}
\label{sec:detect}

The primary comparison retrains YOLO11n, YOLO26n and RT-DETR-l on visible images with
all 300 test images held out of training and validation. YOLO11n is repeated over four
seeds; YOLO26n and RT-DETR-l contribute one clean run each. The same split trains four
modality-agnostic YOLO11n runs. Table~\ref{tab:cleansplit} holds the fused images fixed
and asks whether any one of the fourteen inputs beats the visible source scored in the
colour convention the detector was trained on. The table reports the mean and the full
range over training runs for the four-run columns and the paired image interval for
the single-run columns. Table~\ref{tab:detfam} retains the larger pooled audit: it compares four
evaluators and both source conventions, but its validation split overlaps the scored
images and is therefore an optimistic sensitivity analysis.

\begin{table*}[!t]\centering\small

\caption{Primary independent-test comparison on the 300 LLVIP test images. All 300
images were held out before training every evaluator. YOLO11n visible-only and
modality-agnostic models have four runs; YOLO26n and RT-DETR-l visible-only models have
one run each. The header gives the visible-source mAP50. For four-run columns, the
bracket is the full range of observed run contrasts; for one-run columns, it is the
paired image-bootstrap 95 percent interval. BH correction within each one-run
evaluator leaves all 14 fused inputs below the source at $q<0.002$. Negative values
mean that fusion loses; positive values in the final column identify reversals under
the modality-agnostic evaluator.}
\label{tab:cleansplit}
\setlength{\tabcolsep}{1.5pt}
\begin{tabular*}{\linewidth}{@{\extracolsep{\fill}}l cccc@{}}
\toprule
Input & {\shortstack[c]{YOLO11n VIS\\4 runs (0.8657)}} &
{\shortstack[c]{YOLO26n VIS\\1 run (0.8122)}} &
{\shortstack[c]{RT-DETR-l VIS\\1 run (0.8095)}} &
{\shortstack[c]{YOLO11n agnostic\\4 runs (0.8660)}}\\
\midrule
mean fusion & $-0.060$ $[-0.086,-0.029]$ & $-0.147$ $[-0.188,-0.107]$ & $-0.129$ $[-0.166,-0.090]$ & $+0.023$ $[-0.000,+0.038]$\\
activity-only & $-0.046$ $[-0.062,-0.026]$ & $-0.087$ $[-0.119,-0.055]$ & $-0.070$ $[-0.100,-0.044]$ & $-0.021$ $[-0.038,-0.013]$\\
top-hat prior & $-0.048$ $[-0.074,-0.024]$ & $-0.148$ $[-0.190,-0.104]$ & $-0.123$ $[-0.162,-0.086]$ & $+0.026$ $[+0.006,+0.036]$\\
learned probe & $-0.054$ $[-0.082,-0.028]$ & $-0.142$ $[-0.185,-0.101]$ & $-0.119$ $[-0.159,-0.083]$ & $+0.023$ $[+0.005,+0.034]$\\
top-hat, permuted & $-0.058$ $[-0.088,-0.039]$ & $-0.130$ $[-0.170,-0.092]$ & $-0.110$ $[-0.147,-0.073]$ & $+0.014$ $[-0.011,+0.026]$\\
probe, permuted & $-0.060$ $[-0.095,-0.034]$ & $-0.126$ $[-0.163,-0.090]$ & $-0.108$ $[-0.145,-0.072]$ & $+0.015$ $[-0.010,+0.026]$\\
top-hat, w-shuffled & $-0.062$ $[-0.088,-0.037]$ & $-0.143$ $[-0.183,-0.103]$ & $-0.124$ $[-0.163,-0.087]$ & $+0.015$ $[-0.010,+0.025]$\\
probe, w-shuffled & $-0.065$ $[-0.096,-0.038]$ & $-0.144$ $[-0.186,-0.104]$ & $-0.121$ $[-0.159,-0.084]$ & $+0.015$ $[-0.008,+0.026]$\\
ensemble & $-0.050$ $[-0.080,-0.023]$ & $-0.145$ $[-0.189,-0.102]$ & $-0.110$ $[-0.149,-0.074]$ & $+0.025$ $[+0.008,+0.035]$\\
IFCNN & $-0.076$ $[-0.091,-0.050]$ & $-0.092$ $[-0.132,-0.052]$ & $-0.095$ $[-0.131,-0.059]$ & $+0.016$ $[-0.001,+0.024]$\\
DenseFuse & $-0.052$ $[-0.078,-0.029]$ & $-0.079$ $[-0.118,-0.041]$ & $-0.088$ $[-0.126,-0.053]$ & $+0.013$ $[-0.007,+0.023]$\\
classif. saliency & $-0.061$ $[-0.089,-0.030]$ & $-0.149$ $[-0.189,-0.109]$ & $-0.132$ $[-0.170,-0.094]$ & $+0.023$ $[-0.002,+0.037]$\\
PIAFusion & $-0.104$ $[-0.123,-0.075]$ & $-0.198$ $[-0.244,-0.151]$ & $-0.207$ $[-0.252,-0.163]$ & $+0.029$ $[+0.001,+0.045]$\\
SPGFusion & $-0.136$ $[-0.146,-0.126]$ & $-0.223$ $[-0.269,-0.177]$ & $-0.246$ $[-0.290,-0.204]$ & $-0.009$ $[-0.044,+0.009]$\\
\bottomrule
\end{tabular*}

\end{table*}

\begin{table*}[!t]\centering\small

\caption{The optimistic pooled audit on the same 300 LLVIP images. Its
validation split overlaps the scored images; Table~\ref{tab:cleansplit} is the primary
comparison. The fused rows elsewhere in the audit use a single channel. Printing the
visible source in both conventions exposes the $0.038$ to $0.098$ mAP50 cost of that
choice. ``Below'' counts the fourteen fused inputs whose paired interval lies below the
colour-scored source; the RT-DETR-l count uses the three-run mean. The per-method values
for the two single-run YOLO evaluators are in Table~\ref{tab:dualconv}; RT-DETR-l
intervals and BH-adjusted $q$ values are in Supplementary Table~S23.}\label{tab:detfam}
\begin{tabular*}{\linewidth}{@{\extracolsep{\fill}}l SSSS c@{}}
\toprule
Evaluator & {VIS single} & {VIS colour} & {min deficit} & {max deficit} & {below}\\
\midrule
YOLO11n & 0.8289 & 0.8809 & 0.048 & 0.152 & 14 / 14\\
YOLO26n & 0.8209 & 0.8723 & 0.041 & 0.143 & 14 / 14\\
RT-DETR-l, 3 runs & 0.8624 & 0.9065 & 0.014 & 0.092 & 14 / 14\\
modality-agnostic & 0.7561 & 0.8541 & {--} & {--} & 0 / 14\\
\bottomrule
\end{tabular*}

\end{table*}

\begin{table*}[!t]\centering\small

\caption{Per-input pooled contrasts for the two single-run YOLO evaluators under both
visible-source conventions on the same 300 LLVIP images. The top line of each cell is
$s$, the mean mAP50 difference and its Benjamini--Hochberg $q$ value when every input,
including VIS, is scored in a single channel; the bottom line is $c$, the corresponding
values when VIS is scored in colour. Each convention is a separate family of 42 tests
(14 inputs over the three visible-trained evaluators in the reported run); the displayed
$q$ values therefore use that complete family. The RT-DETR-l three-run summary is in
Table~\ref{tab:detfam}, and its per-input single-run intervals and $q$ values are in
Supplementary Table~S23.}\label{tab:dualconv}
\setlength{\tabcolsep}{3pt}
\begin{tabular*}{\linewidth}{@{\extracolsep{\fill}}l cc@{}}
\toprule
Input & {YOLO11n ($s/c$)} & {YOLO26n ($s/c$)}\\ 
\midrule
mean fusion & \shortstack[l]{$-0.003$ $(0.975)$\\$-0.054$ $(<0.001)$} & \shortstack[l]{$-0.026$ $(0.756)$\\$-0.077$ $(0.003)$}\\
activity-only & \shortstack[l]{$-0.003$ $(0.975)$\\$-0.053$ $(<0.001)$} & \shortstack[l]{$+0.009$ $(0.975)$\\$-0.041$ $(0.016)$}\\
top-hat prior & \shortstack[l]{$+0.000$ $(0.996)$\\$-0.050$ $(<0.001)$} & \shortstack[l]{$-0.005$ $(0.975)$\\$-0.056$ $(0.006)$}\\
learned probe & \shortstack[l]{$-0.005$ $(0.975)$\\$-0.055$ $(<0.001)$} & \shortstack[l]{$-0.009$ $(0.975)$\\$-0.060$ $(0.003)$}\\
top-hat, permuted & \shortstack[l]{$+0.002$ $(0.975)$\\$-0.048$ $(<0.001)$} & \shortstack[l]{$-0.015$ $(0.975)$\\$-0.066$ $(0.003)$}\\
probe, permuted & \shortstack[l]{$+0.001$ $(0.983)$\\$-0.049$ $(<0.001)$} & \shortstack[l]{$-0.019$ $(0.975)$\\$-0.070$ $(0.003)$}\\
top-hat, w-shuffled & \shortstack[l]{$-0.007$ $(0.975)$\\$-0.057$ $(<0.001)$} & \shortstack[l]{$-0.021$ $(0.973)$\\$-0.072$ $(0.003)$}\\
probe, w-shuffled & \shortstack[l]{$-0.005$ $(0.975)$\\$-0.055$ $(<0.001)$} & \shortstack[l]{$-0.028$ $(0.635)$\\$-0.078$ $(<0.001)$}\\
ensemble & \shortstack[l]{$-0.000$ $(0.983)$\\$-0.051$ $(<0.001)$} & \shortstack[l]{$-0.011$ $(0.975)$\\$-0.062$ $(0.003)$}\\
IFCNN & \shortstack[l]{$-0.049$ $(0.063)$\\$-0.099$ $(<0.001)$} & \shortstack[l]{$-0.039$ $(0.364)$\\$-0.090$ $(<0.001)$}\\
DenseFuse & \shortstack[l]{$-0.008$ $(0.975)$\\$-0.059$ $(<0.001)$} & \shortstack[l]{$+0.008$ $(0.975)$\\$-0.043$ $(0.028)$}\\
classif. saliency & \shortstack[l]{$-0.006$ $(0.975)$\\$-0.057$ $(<0.001)$} & \shortstack[l]{$-0.030$ $(0.536)$\\$-0.081$ $(<0.001)$}\\
PIAFusion & \shortstack[l]{$-0.037$ $(0.420)$\\$-0.088$ $(<0.001)$} & \shortstack[l]{$-0.006$ $(0.975)$\\$-0.056$ $(0.008)$}\\
SPGFusion & \shortstack[l]{$-0.101$ $(<0.001)$\\$-0.151$ $(<0.001)$} & \shortstack[l]{$-0.092$ $(<0.001)$\\$-0.143$ $(<0.001)$}\\
\bottomrule
\end{tabular*}

\end{table*}

Figure~\ref{fig:claim} brings the clean-split evidence and the pooled sensitivity audit
together.

\begin{figure*}[!t]\centering

\includegraphics[width=\linewidth]{fig_claim_v2.pdf}
\caption{Clean-split evidence first, then the optimistic protocol audit. (a) The three
clean visible-only evaluators on the 300 held-out images. YOLO11n is the mean over four
runs with the full run range; YOLO26n and RT-DETR-l are single runs with paired image
intervals. All forty-two evaluator-input contrasts lie below the colour-scored visible
source. (b) The clean four-run modality-agnostic evaluator: red intervals identify the
four inputs that gain in every run, while the activity-only control loses in every run.
(c) The pooled audit, with 95 percent image intervals; the two YOLO families are single
runs and RT-DETR-l is the mean of three runs. This panel is a sensitivity analysis
because validation overlaps the test images. Rows are ordered by the deficit under the
pooled visible-trained YOLO11n.}
\label{fig:claim}

\end{figure*}

\begin{sloppypar}
The clean split is the primary result. Under the reference visible-only evaluator,
YOLO11n, no fused input beats the unfused visible image after averaging the 300 test
images: the mean deficits run from $0.046$ to $0.136$, every one of the four runs stays
on the negative side, and the ranking does not change
(Table~\ref{tab:cleansplit}, Fig.~\ref{fig:claim}a). The independently trained
YOLO26n run also places all fourteen below the source, by $0.079$ to $0.223$ mAP50;
every one of those fourteen contrasts has $q<0.002$ in its own BH family. The
clean RT-DETR-l run places all fourteen below by $0.070$ to $0.246$, again with
$q<0.002$ throughout. Across the three visible-only evaluators, therefore, all 42
input comparisons are negative on the independent test set. The largest losses belong
to the published networks and the reported fidelity index ranks SPGFusion first. The same clean split reverses
the sign under the modality-agnostic evaluator for four inputs, by $0.022$ to $0.029$ in
every run, while the activity-only control loses $0.021$ in every run and the remaining
nine inputs are unresolved (Fig.~\ref{fig:claim}b).

The source alignment changes that conclusion by subgroup. We estimated a translational
residual for all 300 LLVIP source pairs and formed quartiles without using annotations or
detector scores. In the best-aligned quartile, mean fusion, the top-hat prior and the
probe gain $+0.024$, $+0.035$ and $+0.032$ mAP50 in all four clean visible-only runs,
whereas SPGFusion gives $-0.021$ with a $[-0.040,+0.005]$ seed range. In the worst-aligned
quartile those four contrasts are $-0.048$, $-0.033$, $-0.043$ and $-0.165$. The
selection co-varies with scene content: the best-aligned quartile has lower visible
contrast and larger annotated boxes than the full set. Residual registration and
composition therefore contribute to the pooled deficit; it is not a claim of invariance
across these strata (Supplementary Table~S8).

The pooled audit is deliberately broader and deliberately optimistic. Fourteen fused
inputs, the ten rules of Table~\ref{tab:detfam} plus four scrambled controls, give
forty-two contrasts over three visible-trained architectures. All forty-two lie below
the colour-scored source (Fig.~\ref{fig:claim}c), with deficits of $0.014$ to $0.152$.
Scoring the visible source in the colour convention raises it from $0.829$ to $0.881$
mAP50 under YOLO11n, and by $0.051$, $0.044$ and $0.098$ under YOLO26n, RT-DETR-l and
the modality-agnostic detector, which enlarges every deficit. Rescoring with the detector
library rather than our pooled implementation preserves the sign, while changing a
deficit by up to $0.015$ mAP50 on our rules and $0.035$ on the two published networks
(Supplementary Secs.~14 and 15). The pooled lower bound, $0.014$, is the activity-only
control under RT-DETR-l; its three runs give $-0.013$, $-0.018$ and $-0.011$, with a
mean of $-0.014$ $[-0.027,-0.001]$. Because that audit selects and scores on the same
300 images, it is not used as the primary estimate; its role is to show that the sign
does not depend on the clean split alone.

A qualification belongs with these numbers. Our own four rules are among the fourteen
that lose, so the negative result does not favour the authors' family: in the clean split
the top-hat prior and probe lose $0.048$ and $0.054$ under four-run YOLO11n,
$0.148$ and $0.142$ under YOLO26n, and $0.123$ and $0.119$ under RT-DETR-l. The
parameter-free activity weighting is in fact the best of the four author rules in the
two clean single-run evaluators, by $0.049$ to $0.061$ over the targetness rules. The
targetness weighting is therefore not the best downstream member of its own family.

At the stricter mAP50-95 criterion the sign is evaluator-specific. Fused inputs under
the clean YOLO11n reference are ahead of the visible source at the reported seed by
about $0.03$ for many rules, whereas every fused input is below the source under clean
YOLO26n (Supplementary Sec.~11). We therefore do not generalize the headline across
criteria; the primary claim is stated at mAP50.
\end{sloppypar}

\begin{sloppypar}
\paragraph{Where the deficit is.} The stratification in Fig.~\ref{fig:regime} is a
pooled sensitivity analysis, not a second primary result. Its unadjusted counts reach
fourteen in the high-contrast, crowded and small-target strata, while the corresponding
Benjamini--Hochberg counts in Supplementary Table~S16 are fourteen in the
high-contrast and small-target strata and five to fourteen in the crowded stratum.
The clean split does not reproduce one invariant partition. In the order low contrast,
high contrast, one pedestrian, three or more, small and large targets, the number of
inputs resolved below is $1,1,0,3,2,0$ under clean YOLO11n, $13,14,11,14,14,14$ under
clean YOLO26n and $2,14,2,14,14,5$ under clean RT-DETR-l; no input is resolved above
in any clean stratum (Supplementary Table~S29). The aggregate negative sign is
therefore robust across evaluators, but which subgroup supplies the resolution is not.
A continuous pooled split of the 300 images into six contrast bins shows the same lack
of a monotone contrast curve, and the sign of the within-sequence association also
changes by sequence (Supplementary Tables S18 and S19). The strata use visible RMS
contrast boundaries $0.0830$ and $0.1953$, a small-target cutoff of $26{,}902.72$
pixels and a crowded threshold of three or more pedestrians. The segmentation
repetition (Fig.~\ref{fig:regime}c) is likewise retained as a descriptive control.
\end{sloppypar}

\begin{figure*}[!t]\centering

\includegraphics[width=\linewidth]{fig_regime_v2.pdf}
\caption{Where the deficit is in the pooled sensitivity audit. (a) Each dot is one of
the fourteen fused inputs: green, resolved below the colour-scored visible source;
grey, unresolved; orange, resolved above. The matrix shows every regime and every
visible-trained architecture. Regimes are computed from the annotation and unfused
visible source before any contrast; counts are in Supplementary Table~S16. The panel
uses uncorrected 95 percent intervals, whereas Table~S16 gives BH-adjusted counts for
the overlapping strata. These runs train or select on the scored images and are not the
primary independent-test result; clean counterparts are in Supplementary Table~S29.
(b) The 300 test images in six bins of fifty by visible RMS contrast. Grey lines show
all fourteen inputs; mean fusion, the learned probe and SPGFusion are highlighted, with
the shaded band giving the SPGFusion interval. The bins are also sequences, so this is
a stratification rather than a dose-response. (c) Segmentation on 361 MSRS-test
images, mean fusion against the unfused visible source, by visible-contrast tercile.}
\label{fig:regime}

\end{figure*}

\begin{sloppypar}
\paragraph{What the seed dimension changes.} The clean split removes the optimistic
training overlap, and repeating each detector four times shows which signs are robust.
For the clean-split modality-agnostic detector, four rules exceed the colour-scored
source by $0.022$ to $0.029$ in every run, so the reversal survives the clean split,
while the activity-only control loses in every run (Fig.~\ref{fig:robust}a,
Supplementary Table~S21). This is a narrower reversal than the pooled detector gives:
ten inputs clear the source there, but only four survive full test isolation. The
RT-DETR-l family, which carried the one contrast of the pooled negative result
whose interval included zero, was retrained twice more: the control sits below the
source in all three runs, by $0.013$, $0.018$ and $0.011$ mAP50, and the mean over runs
excludes zero at $-0.014$ $[-0.027,-0.001]$, so all forty-two contrasts are resolved
(Fig.~\ref{fig:robust}b, Supplementary Table~S22). This count belongs to the pooled
audit, not to the primary clean split. The deficits of that architecture
remain the mildest of the three, and the qualification that belongs with the row is
that the colour-scored visible source itself moves by $0.020$ mAP50 between runs, which
is larger than the control's deficit, so a resolved interval on $0.014$ is a statement
at the edge of the resolution this benchmark offers.
\end{sloppypar}

\begin{figure*}[!t]\centering

\includegraphics[width=\linewidth]{fig_robust_v2.pdf}
\caption{What is stable across seeds, runs and index constructions. (a) Clean-split
modality-agnostic detector: paired difference against the colour-scored visible source,
mean over four seeds, with the two-level interval that resamples seeds and then images.
Red intervals mark the four inputs that remain above zero. (b) RT-DETR-l over three
training runs of the same pooled recipe; this panel is a sensitivity analysis because
those validation splits overlap the test images. (c) Rank positions of the same eight
fused variants under the reported index, its noise-calibration departure alone, and the
reference pixel VIF. The red path shows SPGFusion moving from first to last; the two
further departures that leave the order unchanged are in Supplementary Table~S6.}
\label{fig:robust}

\end{figure*}

\begin{sloppypar}
\paragraph{The second benchmark.} The negative result replicates on FLIR-ADAS aligned,
1500 training and 300 test pairs, single class, at native $640\times512$
(Table~\ref{tab:flir}). Under both visible-trained architectures the unfused visible
source is again the best of the twelve inputs and every fused input lies below it;
under YOLO11n the tabulated mAP50 differences run from $-0.0498$ for the probe to
$-0.0928$ for SPGFusion; the paired bootstrap estimates for the same endpoints are
$-0.0515$ $[-0.0840,-0.0161]$ and $-0.0928$ $[-0.1251,-0.0627]$, respectively. The
SPGFusion estimate is almost exactly the $0.0927$ it loses on the first benchmark.
Under YOLO26n the point-estimate differences run from $-0.0084$ for the probe to
$-0.0554$ for the classification-saliency rule; intervals were not computed for this
evaluator. The two visible-trained architectures do not agree on the internal order,
since PIAFusion is tenth of twelve under one and sixth under the other, and what they
agree on is the conclusion about fusion itself. Under the
modality-agnostic detector of that benchmark the infrared source becomes the best input
by a wide margin ($0.763$ against $0.395$) and all fused inputs gain on the visible
source, by $+0.141$ to $+0.256$, so the reversal replicates as well. Its detectors were
trained with a different budget and their weights are no longer available, so the
second benchmark is reported as measured and uncorrected. Table~\ref{tab:cross}
collects the five contrasts that carry the argument under all four combinations of
benchmark and evaluator: every one of them changes sign at least once.
Table~\ref{tab:flir} is in the single-channel convention. The FLIR checkpoints are no
longer available, so the second benchmark cannot be colour-rescored; it therefore
replicates the sign of the primary result and not its colour-convention magnitude.

\begin{table*}[!t]\centering\small

\caption{Replication on FLIR-ADAS aligned, 1500 training and 300 test pairs,
single class person, native $640\times512$. Three detectors are trained on this
benchmark with the same recipe as the first, two of them on the visible modality alone
and differing in architecture. Entries are mAP50 point estimates; best in bold. Signed
contrasts quoted in the text are fusion minus the VIS source in the same column.
Bootstrap estimates in Table~\ref{tab:cross} can differ from differences of rounded
table entries by up to $0.003$.}\label{tab:flir}
\begin{tabular}{@{}l ccc@{}}
\toprule
Input & {YOLO11n} & {YOLO26n} & {agnostic}\\
\midrule
VIS source & \textbf{0.4686} & \textbf{0.3769} & 0.3950\\
IR source & 0.3382 & 0.2955 & \textbf{0.7630}\\
mean fused & 0.3834 & 0.3424 & 0.6511\\
activity-only fused & 0.4018 & 0.3541 & 0.6162\\
top-hat fused & 0.4156 & 0.3571 & 0.6305\\
probe fused & 0.4188 & 0.3686 & 0.6255\\
ensemble fused & 0.4156 & 0.3666 & 0.6296\\
IFCNN fused & 0.4062 & 0.3449 & 0.6217\\
DenseFuse fused & 0.4131 & 0.3636 & 0.6486\\
CSF fused & 0.3813 & 0.3215 & 0.6463\\
PIAFusion fused & 0.3773 & 0.3544 & 0.5528\\
SPGFusion fused & 0.3758 & 0.3459 & 0.5359\\
\bottomrule
\end{tabular}

\end{table*}

\begin{table*}[!t]\centering\small

\caption{The same five contrasts under all four combinations of detection
benchmark and evaluator. A star marks a paired interval that excludes zero; each cell
is the mean of the bootstrap distribution of that contrast on the same estimator as its
interval, which can differ from the difference of the two pooled table values by up to
$0.003$. Every contrast changes sign at least once.}\label{tab:cross}
\begin{tabular*}{\linewidth}{@{\extracolsep{\fill}}l cccc@{}}
\toprule
Contrast & {LLVIP, vis} & {LLVIP, agn} & {FLIR, vis} & {FLIR, agn}\\
\midrule
SPGFusion $-$ VIS & $-0.0927$* & $+0.0924$* & $-0.0928$* & $+0.1404$*\\
mean $-$ VIS & $-0.0023$ & $+0.1330$* & $-0.0862$* & $+0.2556$*\\
probe $-$ mean & $-0.0017$ & $+0.0031$ & $+0.0347$* & $-0.0265$*\\
probe $-$ activity-only & $-0.0027$ & $+0.0388$* & $+0.0172$* & $+0.0089$\\
top-hat $-$ mean & $+0.0033$ & $+0.0003$ & $+0.0316$* & $-0.0211$*\\
\bottomrule
\end{tabular*}

\end{table*}
\end{sloppypar}

\section{Sources of Evaluator and Index Dependence}
\label{sec:ranking}

\begin{sloppypar}
The evaluator and the index are separable, and the same experiment design measures
both, because the fused images are held fixed in every comparison of this section.

\paragraph{The evaluator.} Three factors move the ranking of the same sixteen inputs,
fourteen fused rules and the two sources. The training modality moves
it most. In the primary clean split, four of the fourteen fused inputs gain over the
visible source by $0.022$ to $0.029$ in every run, while the activity-only control loses
in every run (Fig.~\ref{fig:claim}b). The pooled modality-agnostic network gives the
larger but optimistic estimate: ten inputs gain by $0.023$ to $0.040$, and the infrared
source becomes the best input of all ($0.9415$ against $0.8953$ for the best fused
input), although it is the worst under the visible-trained one. The clean split is
therefore the number we retain; the pooled result only shows the direction before the
overlap is removed. In the pooled sensitivity audit, the architecture moves the
ranking less but measurably: the
probe-against-mean contrast is $+0.005$ $[-0.002,+0.014]$ under YOLO11n and $+0.021$
$[+0.014,+0.027]$ under pooled YOLO26n, whose four-seed ranges are disjoint, so the ordering
inside the fused group changes with the architecture by more than that architecture's
own seed spread, though the margin before rounding is $0.0005$ mAP50. In that pooled
audit the deficits under RT-DETR-l are the mildest of the three architectures,
SPGFusion by $0.091$ against $0.152$ under YOLO11n, and the activity-only control that
is mid-pack elsewhere becomes the best fused input under RT-DETR-l in the
single-channel convention ($0.8836$ against $0.8624$, a gain of $0.021$), while in the
colour convention its three-run mean is $0.014$ below the source. The independent-test
contrasts are in Table~\ref{tab:cleansplit}; because only YOLO11n is repeated there,
we do not promote the pooled architecture contrast above a sensitivity analysis. The
clean 14-input rankings correlate at $\rho=0.45$ between YOLO11n and YOLO26n, $0.60$
between YOLO11n and RT-DETR-l, and $0.94$ between YOLO26n and RT-DETR-l, so architecture
changes the magnitude and part of the ordering without breaking the negative sign.

\paragraph{The index.} The second headline is about the metric rather than the network.
We rescored all fourteen LLVIP inputs at $512$ px under the reported index and the
reference construction. The endpoints reverse: SPGFusion is first under the reported
index ($0.812$) and last under reference VIF ($0.170$), while PIAFusion moves from second
($0.706$) to first ($0.372$). The same reversal holds on the three full fidelity
benchmarks. SPGFusion is first under the reported index on TNO, MSRS and RoadScene
($0.785$, $0.831$ and $0.829$) and last under reference VIF
($0.161$, $0.231$ and $0.217$); PIAFusion moves from second to first on TNO and MSRS
(Supplementary Table~S7).
Two of the three departures of the reported index from the reference construction leave
the ordering unchanged, and the third accounts for the reversal: holding the numeric
noise variance at $2.0$ while the inputs are in $[0,1]$ collapses the spread. The
counterexample therefore belongs to the reported index rather than to fidelity metrics in
general, which is why the first headline is stated for the index the article reports.

The architecture effect is small next to the seed effect in the pooled audit, which is
why we quantify both with the same estimator there. Table~\ref{tab:seed} gives six
contrasts under the two YOLO
families, each over four independently trained runs of the same recipe, so that the
range in a cell is the spread the detector introduces by itself. Five of the six
contrasts agree in sign and overlap in range, including whether the probe reaches the
unfused source ($-0.022$ against $-0.017$) and whether it beats the plain activity rule
($-0.017$ against $-0.011$). One separates: the probe against mean fusion is $+0.005$
$[-0.002,+0.014]$ under YOLO11n and $+0.021$ $[+0.014,+0.027]$ under YOLO26n, disjoint
ranges, with a margin before rounding of $0.0005$ mAP50. The claim we make from the
architecture half of Sec.~\ref{sec:ranking} is therefore the narrow one the data
support: the ranking inside the fused group moves with the architecture by more than
that architecture's own seed spread for exactly one of six contrasts, and the claim
that matters is the larger one, that the modality of the evaluator decides whether a
fused input is better than the visible source at all.

That table is a pooled sensitivity analysis because its validation split overlaps the
scored images. The independent-test table makes the same architecture comparison on
clean checkpoints, but only YOLO11n there has four runs; YOLO26n and RT-DETR-l have one
run each. We therefore report their clean contrast levels in
Table~\ref{tab:cleansplit} and do not claim that a difference between their internal
rankings is separated from single-run variation.

\begin{table*}[!t]\centering\small

\caption{The architecture effect against the seed effect in the pooled
sensitivity audit, in the single-channel convention of the reported protocol on the
300 LLVIP images. The validation split of these runs overlaps the scored images.
YOLO11n and YOLO26n are each four independently trained runs on the same 1,500 visible
images with the same recipe. Each cell lists seeds 0/1/2/3, then the four-seed mean and
full range. The independent-test counterparts are in Table~\ref{tab:cleansplit}.}
\label{tab:seed}
\setlength{\tabcolsep}{3pt}
\begin{tabular*}{\linewidth}{@{\extracolsep{\fill}}l cc@{}}
\toprule
Contrast & {\shortstack[l]{YOLO11n, seeds 0/1/2/3\\then mean [range]}} & {\shortstack[l]{YOLO26n, seeds 0/1/2/3\\then mean [range]}}\\
\midrule
SPGFusion $-$ VIS & \shortstack[l]{$-0.091/-0.066/-0.136/-0.152$\\$-0.111$ $[-0.152,-0.066]$} & \shortstack[l]{$-0.084/-0.062/-0.035/-0.028$\\$-0.052$ $[-0.084,-0.028]$}\\
mean $-$ VIS & \shortstack[l]{$-0.002/-0.014/-0.057/-0.036$\\$-0.027$ $[-0.057,-0.002]$} & \shortstack[l]{$-0.025/-0.035/-0.055/-0.038$\\$-0.038$ $[-0.055,-0.025]$}\\
probe $-$ VIS & \shortstack[l]{$-0.003/-0.013/-0.043/-0.031$\\$-0.022$ $[-0.043,-0.003]$} & \shortstack[l]{$-0.008/-0.021/-0.028/-0.012$\\$-0.017$ $[-0.028,-0.008]$}\\
probe $-$ mean & \shortstack[l]{$-0.002/+0.001/+0.014/+0.005$\\$+0.005$ $[-0.002,+0.014]$} & \shortstack[l]{$+0.016/+0.014/+0.027/+0.026$\\$+0.021$ $[+0.014,+0.027]$}\\
top-hat $-$ mean & \shortstack[l]{$+0.003/-0.002/+0.025/-0.001$\\$+0.006$ $[-0.002,+0.025]$} & \shortstack[l]{$+0.021/+0.016/+0.034/+0.019$\\$+0.022$ $[+0.016,+0.034]$}\\
probe $-$ activity-only & \shortstack[l]{$-0.002/-0.013/-0.037/-0.015$\\$-0.017$ $[-0.037,-0.002]$} & \shortstack[l]{$-0.018/-0.021/-0.021/+0.017$\\$-0.011$ $[-0.021,+0.017]$}\\
\bottomrule
\end{tabular*}

\end{table*}
\end{sloppypar}

\section{Robustness and Generalization}
\label{sec:additional}
\label{sec:extra}

\begin{sloppypar}
Five further experiments test the two claims from the other side: whether a fidelity
score can predict the downstream ranking at all, whether the deficits depend on the
criterion and on the operating point, how they survive multiplicity correction, and
whether the rule's own sharpness drives them. All of them use the same frozen runs and
the same per-image statistics as the tables above. Unless an experiment explicitly
uses the independent split, these are sensitivity analyses on the pooled checkpoints;
the primary evaluator comparison remains Table~\ref{tab:cleansplit}.

\paragraph{Does fidelity predict the ranking?} The analysis correlates the fidelity of
each fused input with its paired contrast against the colour-scored visible
source. The predictive correlation uses the source-facing VIF, so that the score and the
target are measured against the same reference; Table~\ref{tab:fidfull} also prints the
reported index, the mean of the two source-facing scores. The original fidelity export
predates CSF, so the source-facing scatter uses the other thirteen inputs; CSF is present
in the complete rescoring of Table~\ref{tab:fidfull}. Across those thirteen inputs, the
rank correlation of source-facing reported VIF with that contrast is
$-0.52$ under YOLO11n, $-0.40$ under RT-DETR-l, and $+0.07$ under YOLO26n and under the
modality-agnostic detector. The sign is therefore not stable across evaluators, and the
magnitude never exceeds $0.52$. Within the authors' own family the same correlation is
$-0.10$, $+0.10$, $+0.10$ and $+0.50$ across the four evaluators; within the four
published networks it is $-0.80$, $-0.80$, $-1.00$ and $-0.80$. The mismatch therefore
is not a property of one rule family.

The full fourteen-input rescoring sharpens the role of the construction. Using the
reported score, the rank correlations are $-0.37$ under pooled YOLO11n, $-0.18$ under
RT-DETR-l, $+0.16$ under YOLO26n and $-0.09$ under the pooled modality-agnostic detector.
The reference construction moves them to $-0.04$, $-0.02$, $+0.47$ and $+0.63$. On the
clean evaluators, the reported/reference associations are $-0.13/+0.18$ for four-run
YOLO11n, $-0.11/-0.21$ for YOLO26n, $-0.00/-0.20$ for RT-DETR-l and $+0.07/+0.59$
for the four-run modality-agnostic detector. Thus the reported construction does not
transfer, and the corrected pixel VIF is informative only when the evaluator is
modality-agnostic. SSIM against the visible
source contributes $+0.20$, $+0.46$, $+0.62$ and $+0.19$, and entropy remains between
$-0.37$ and $+0.06$. Fig.~\ref{fig:corr} shows the source-facing scatter.
These are descriptive Spearman rank correlations over a small, deliberately chosen set
of fusion rules, not inferential estimates about a population of methods.

\begin{table*}[!t]\centering\small

\caption{Fidelity and downstream contrast for every fused input. The reported
index and the reference pixel VIF are both computed on all 300 LLVIP test images at
$512$ px and on the same visible-source grid. Detection columns are paired mAP50
contrasts against the colour-scored visible source: the pooled YOLO11n sensitivity
audit, the four-run YOLO11n clean mean, and the clean YOLO26n and RT-DETR-l runs.}
\label{tab:fidfull}
\setlength{\tabcolsep}{4pt}
\setlength{\tabcolsep}{2.5pt}
\begin{tabular*}{\linewidth}{@{\extracolsep{\fill}}l cccccc@{}}
\toprule
Input & {reported VIF} & {reference VIF} & {$\Delta$ pooled} & {$\Delta$ Y11 clean} &
{$\Delta$ Y26 clean} & {$\Delta$ RTD clean}\\
\midrule
mean fusion        & 0.284 & 0.269 & $-0.054$ & $-0.060$ & $-0.147$ & $-0.129$\\
activity-only      & 0.328 & 0.251 & $-0.053$ & $-0.046$ & $-0.087$ & $-0.070$\\
top-hat prior      & 0.343 & 0.275 & $-0.050$ & $-0.048$ & $-0.148$ & $-0.123$\\
learned probe      & 0.358 & 0.272 & $-0.055$ & $-0.054$ & $-0.142$ & $-0.119$\\
top-hat, permuted  & 0.286 & 0.267 & $-0.048$ & $-0.058$ & $-0.130$ & $-0.110$\\
probe, permuted    & 0.287 & 0.267 & $-0.049$ & $-0.060$ & $-0.126$ & $-0.108$\\
top-hat, w-shuffled& 0.286 & 0.270 & $-0.058$ & $-0.062$ & $-0.143$ & $-0.124$\\
probe, w-shuffled  & 0.285 & 0.270 & $-0.055$ & $-0.065$ & $-0.144$ & $-0.121$\\
ensemble           & 0.350 & 0.271 & $-0.051$ & $-0.050$ & $-0.145$ & $-0.110$\\
IFCNN              & 0.379 & 0.259 & $-0.099$ & $-0.076$ & $-0.092$ & $-0.095$\\
DenseFuse          & 0.299 & 0.277 & $-0.059$ & $-0.052$ & $-0.079$ & $-0.088$\\
classif. saliency  & 0.281 & 0.272 & $-0.057$ & $-0.061$ & $-0.149$ & $-0.132$\\
PIAFusion          & 0.706 & 0.372 & $-0.088$ & $-0.104$ & $-0.198$ & $-0.207$\\
SPGFusion          & 0.812 & 0.170 & $-0.152$ & $-0.136$ & $-0.223$ & $-0.246$\\
\bottomrule
\end{tabular*}

\end{table*}

\begin{figure*}[!t]\centering

\includegraphics[width=\linewidth]{fig_corr_v2.pdf}
\caption{Fidelity does not transfer across evaluators. (a) Source-facing reported-index
VIF against the paired difference from the colour-scored visible source on the thirteen
inputs in the original fidelity export. CSF is the one omitted input, not a dropped row.
(b) Spearman rank correlations for VIF, SSIM against the visible source and entropy,
using the evaluator symbols defined in (a). SSIM stays positive under all four
evaluators, with a threefold range; the reported VIF is negative under YOLO11n and
RT-DETR-l and positive under the other two, and entropy stays between $-0.37$ and
$+0.06$.
Table~\ref{tab:fidfull} reports both index constructions on the complete fourteen-input,
300-image table.}
\label{fig:corr}

\end{figure*}

\paragraph{The criterion.} The criterion changes the ordering rather than establishing a
second sign law. Under pooled YOLO11n SPGFusion loses $0.151$ mAP at IoU 0.50 and
$0.163$ at 0.55 but is level with the source beyond IoU 0.85, and under the
modality-agnostic detector the same rule goes from $-0.013$ at IoU 0.50 to $+0.190$ at
0.80 (Table~\ref{tab:criterion}). Fusion therefore changes which pedestrians are found
at a loose criterion rather than how tightly the found ones are localised, and the
advantage of the modality-agnostic evaluator is largest where localisation is hardest
(Fig.~\ref{fig:criterion}). On the clean runs every fused point estimate is negative at
mAP50-95 under YOLO26n and RT-DETR-l, while several are positive under the YOLO11n
seed; the criterion effect is therefore architecture-specific and is not used as a
headline.

\begin{table*}[!t]\centering\small

\caption{The criterion sweep. Columns 2 to 4 give the visible source at IoU
0.50, 0.75 and averaged over 0.50 to 0.95; columns 5 to 7 give the contrast of
SPGFusion against that source at IoU 0.50, 0.75 and 0.85. mAP50 and mAP, so the
columns are comparable within a row only through the contrast.}\label{tab:criterion}
\setlength{\tabcolsep}{4pt}
\begin{tabular*}{\linewidth}{@{\extracolsep{\fill}}l SSSS SSS@{}}
\toprule
Evaluator & {vis 0.50} & {vis 0.75} & {vis 0.50-95} & {$\Delta$ 0.50} & {$\Delta$ 0.75} & {$\Delta$ 0.85}\\
\midrule
YOLO11n & 0.8811 & 0.5064 & 0.4893 & -0.151 & -0.072 & +0.001\\
YOLO26n & 0.8727 & 0.4498 & 0.4587 & -0.141 & -0.056 & -0.009\\
RT-DETR-l & 0.9065 & 0.4723 & 0.4836 & -0.096 & +0.037 & +0.071\\
modality-agnostic & 0.8550 & 0.4686 & 0.4533 & -0.013 & +0.116 & +0.137\\
\bottomrule
\end{tabular*}

\end{table*}

\begin{figure*}[!t]\centering

\includegraphics[width=\linewidth]{fig_criterion_v2.pdf}
\caption{The deficit depends on the criterion. Paired difference against the
colour-scored visible source as a function of the IoU criterion at which the average
precision is taken, for four representative rules, under (a) the primary
visible-trained detector and (b) the modality-agnostic detector of the main tables.
Fusion loses detections at a loose criterion, where the visible source is already
strong, and becomes level with or better than it where tight localisation is required
and every input scores lower.}
\label{fig:criterion}

\end{figure*}

\paragraph{Where the deficit comes from.} At the maximum-F1 operating point of each
input at IoU 0.50, the loss is dominated by missed pedestrians rather than by false
positives. The visible source reaches $0.827$ recall and $0.908$ precision under
YOLO11n, missing $129$ of the $744$ annotated pedestrians; SPGFusion reaches $0.642$
and $0.846$ and misses $266$, IFCNN $0.695$ and $0.817$ with $227$ missed, mean fusion
$0.708$ and $0.886$ with $217$, and the learned probe $0.741$ and $0.849$ with $193$.
The recall loss is therefore four to seven times the precision loss for every fused
input, and larger than the whole contrast in mAP50, which is consistent with the
criterion sweep: a fused image suppresses pedestrians that the visible source finds.
Under the modality-agnostic detector the same four inputs gain recall instead
($+0.023$ to $+0.042$) while losing a little precision (Fig.~\ref{fig:operating}).

\begin{figure*}[!t]\centering

\includegraphics[width=\linewidth]{fig_operating_v2.pdf}
\caption{Where the deficit comes from, at the maximum-F1 operating point of each input
at IoU 0.50: (a) recall and (b) precision, for the unfused visible source and four
fused inputs, under the four evaluators. Fusion costs recall, not precision, under
every visible-trained architecture, and gains recall under the modality-agnostic one.}
\label{fig:operating}

\end{figure*}

\paragraph{Multiplicity.} The pooled sensitivity family is the forty-two contrasts over
the three visible-trained architectures. We correct the single-channel
and colour-scored comparisons as two separate 42-test families, using the
Benjamini--Hochberg procedure on bootstrap two-sided $p$ values from the same 1000
paired image draws, at a false-discovery rate of $0.05$. In the single-channel family,
only three contrasts reach $q<0.05$, all three for SPGFusion; the apparent gains of
several other fused inputs in Table~\ref{tab:dualconv} are unresolved. In the
colour-scored family, forty of the forty-two have $q<0.05$ and thirty-one have
$q<0.01$. The largest $q$ belongs to the activity-only control under the reported
RT-DETR-l run ($q=0.210$), which is the contrast that the three-run repetition of
Supplementary Sec.~19 resolves, and the next three are $0.055$, $0.048$ and $0.045$
for the probe-permuted control, the ensemble and the top-hat prior under that same
architecture. The number of resolved colour-scored contrasts is also stable in the
number of bootstrap draws: $41$, $41$, $41$ and $40$ of the forty-two are resolved as
the resampling is truncated to $100$, $200$, $400$ and $1000$ draws. The correction
therefore changes which single-channel wins we regard as established, but not the
colour-scored conclusion.

\paragraph{The sharpness of the rule.} The temperature $\tau$ of Eq.~\eqref{eq:rule}
sets how near-binary the per-pixel weighting is, and the two families of rule respond to
it very differently. Over $\tau\in\{3,1,0.3,0.1\}$, that is a thirty-fold range, the
top-hat prior stays within $0.002$ mAP50 of its own best under the visible-trained
detector and within $0.014$ under the modality-agnostic one, and the learned probe
within $0.005$ and $0.012$; the plain activity rule, which is the same formula with
$S\equiv1$, loses up to $0.045$ and $0.089$ over the same range
(Fig.~\ref{fig:tau}). What the targetness map contributes is therefore the stability of
the weighting rather than a fixed gain: at the reported temperature the three rules are
within a few thousandths of each other, and they separate only when the rule is pushed
towards a hard decision. The underlying weight distributions make the mechanism explicit.
At $\tau=3,1,0.3,0.1$, the per-image mean $|w_v-0.5|$ and its cross-image standard
deviation are $0.012\pm0.001$, $0.037\pm0.004$, $0.117\pm0.013$ and $0.268\pm0.023$ for
activity-only, $0.005\pm0.002$, $0.015\pm0.005$, $0.047\pm0.016$ and $0.125\pm0.038$ for
the top-hat prior, and $0.003\pm0.001$, $0.008\pm0.003$, $0.026\pm0.011$ and
$0.073\pm0.029$ for the probe on TNO. On MSRS the corresponding means are
$0.018,0.053,0.159,0.316$; $0.003,0.010,0.034,0.092$; and $0.002,0.006,0.019,0.049$.
The spread therefore separates the three rules even before detection is scored, which is
why the comparison is repeated at matched deviation in Supplementary Table~S2.

\begin{figure*}[!t]\centering

\includegraphics[width=\linewidth]{fig_tau_v2.pdf}
\caption{The rule's own sharpness. Change in mAP50 relative to mean fusion as the
temperature of Eq.~\eqref{eq:rule} is varied over a thirty-fold range on a logarithmic
axis, under (a) the visible-trained and (b) the modality-agnostic evaluator. The zero
line is mean fusion, which has no temperature; the top-hat prior and the learned probe
hold their value across the whole range, while the plain activity rule, the same formula
with $S\equiv1$, degrades sharply once the weighting becomes hard.}
\label{fig:tau}

\end{figure*}

\paragraph{A larger visible-trained backbone.} The negative result could be a property
of small detectors, so the family recipe of Sec.~\ref{sec:method} was repeated with a
larger backbone: YOLO11s, trained for the same 40 epochs at the same batch size and
image size on the same 1500 visible images. The larger model does not raise the ceiling
on this training set, because the colour-scored visible source scores $0.8719$ mAP50
against $0.8811$ for YOLO11n. Every one of the fifteen fused inputs still lies below
the source in point estimate, SPGFusion by $0.071$ and IFCNN by $0.040$ with intervals
that exclude zero, and the infrared source loses $0.115$; but only four of the fifteen
contrasts resolve, against fifteen of fifteen under the reported architecture, because
the larger backbone shrinks the deficits, to between $0.011$ and $0.071$ mAP50. The
result therefore persists in sign and loses resolution rather than reversing, and the
size of the penalty is a property of the detector family as much as of the fused image
(Table~\ref{tab:scale}, Fig.~\ref{fig:scale}a).

\paragraph{The inference resolution.} The comparison above is at native test
resolution with the detector run at its training size, $640$\,px. Repeating the whole
scoring at $512$ and $960$\,px changes the level of every input without changing the
conclusion: the visible source reads $0.8636$, $0.8811$ and $0.7582$ mAP50 at the
three sizes, that is, the detector is tuned for its training size and degrades at both
others, and all fifteen fused inputs are resolved below the source at all three sizes,
with SPGFusion losing $0.140$, $0.151$ and $0.234$ (Table~\ref{tab:scale},
Fig.~\ref{fig:scale}b). The deficits grow as the evaluation moves away from the
detector's operating point while the sign does not move at all, so the negative result
is not an artefact of one evaluation size.

\begin{table*}[!t]\centering\small

\caption{The scale experiments. Upper block: the reported architecture and the
larger backbone, on the same 1500 visible training images and the same test images,
with the colour-scored visible source of each and the contrast of four representative
inputs. Lower block: the primary visible-trained detector scored at three inference
sizes, with the same source and contrasts. A star marks a contrast whose interval
excludes zero; the reported run uses 1000 draws and the resolution sweep 400.
}\label{tab:scale}
\setlength{\tabcolsep}{4pt}
\begin{tabular*}{\linewidth}{@{\extracolsep{\fill}}l SSSS S@{}}
\toprule
Setting & {VIS source} & {SPGFusion} & {IFCNN} & {IR source} & {resolved of 15}\\
\midrule
YOLO11n, reported run & 0.8811 & -0.151* & -0.099* & -0.174* & 15\\
YOLO11s, larger backbone & 0.8719 & -0.071* & -0.040* & -0.115* & 4\\
\addlinespace
YOLO11n at imgsz 512 & 0.8636 & -0.140* & -0.119* & -0.211* & 15\\
YOLO11n at imgsz 640 & 0.8811 & -0.151* & -0.099* & -0.174* & 15\\
YOLO11n at imgsz 960 & 0.7582 & -0.234* & -0.144* & -0.199* & 15\\
\bottomrule
\end{tabular*}

\end{table*}

\begin{figure*}[!t]\centering

\includegraphics[width=\linewidth]{fig_scale_v2.pdf}
\caption{The two scale experiments. (a) Every fused input under the larger
visible-trained backbone (YOLO11s, mean over the run with its interval) against the same
input under the reported architecture; rows are ordered as in Fig.~\ref{fig:claim}. The
larger backbone keeps every deficit in sign and shrinks most of them below the
resolution of the comparison. (b) The same inputs at three inference sizes with the
primary visible-trained detector.}
\label{fig:scale}

\end{figure*}

\paragraph{Sequence by sequence.} The stratification of Sec.~\ref{sec:detect} pools the
test set into terciles; the benchmark is also organised into sequences, and
Supplementary Tables S18 and S19 resolve the same contrast inside each of the seven LLVIP
sequences that contribute at least twenty test images. The visible source itself scores between
$0.690$ and $0.954$ mAP50 across them, and the sign of the fused contrast follows the
sequence rather than the pooled average: in the sequence where the visible source is
weakest ($0.690$) every fused input gains, by $+0.059$ to $+0.172$ mAP50, whereas two
sequences lose between $0.19$ and $0.37$. The sequence with the largest deficit is not
the darkest one but the one where the source is mid-ranked ($0.815$), which is the same
statement the contrast bins make: the pooled negative is an average over a benchmark in
which individual sequences behave in opposite directions.

For all four families the per-sequence sign is the same: six of the seven sequences are
negative and the one positive sequence is the one with the weakest visible source
($0.690$). The five strata of Fig.~\ref{fig:regime} were not pre-registered. We report
them as a description of the test set, report the BH-adjusted counts within each
architecture in Supplementary Sec.~16, and retain the larger small-target and
high-contrast effects as hypotheses for a pre-registered study rather than causal
estimates of a regime effect.

\paragraph{The segmentation protocol.} The third task uses a lightweight U-Net with one
input channel and eight foreground classes plus background, trained on the MSRS training
pairs and evaluated on the 361 MSRS-test pairs against their segmentation labels. The
main table reports the strict convention that averages over all nine classes and counts a
class absent from both prediction and label as zero, which is why absolute mIoU is about
$0.12$. The usual present-classes convention gives about $0.30$ and reverses the visible
and infrared ranking, so the segmentation numbers are read as a third ranking under a
stated convention (Supplementary Sec.~17).

\paragraph{How large a deficit is resolvable here.} Because every interval in the article
comes from resampling the same 300 images, the size of a contrast that the benchmark can
resolve is measurable rather than assumed. Over the sixty contrasts of the study the
paired bootstrap standard deviation has a median of $0.0162$ mAP50, a minimum of
$0.0105$ and a maximum of $0.0240$, so the 95 percent detection floor of a single
contrast has a median of $0.032$ mAP50 and lies between $0.021$ and $0.047$ depending
on the input and on the evaluator. This is the bound that the reported deficits have to
beat: the smallest resolved deficit of the article, $0.022$ under a clean split, is
resolvable only because its own standard deviation is at the low end of that range, and
deficits near $0.01$ are below the floor of this benchmark at 300 images rather than
absent. SPGFusion's $0.15$ mAP50 under the visible-trained detector is nine standard
deviations, while the activity-only control's $0.014$ under RT-DETR-l is about one,
which is why that contrast needed three runs.

\paragraph{Two evaluators disagree image by image.} A pooled contrast hides what
happens per image. Scoring every test image by its own maximum-F1 point and repeating
the SPGFusion contrast image by image, two evaluators agree on the sign of the contrast
in only $57$ to $69$ percent of the images, and agree that the fused input is worse in
only $8$ to $26$ percent. The per-image contrasts of the three visible-trained
architectures correlate between $+0.34$ and $+0.51$, and those of the modality-agnostic
detector with any visible-trained one between $+0.04$ and $+0.26$. The negative result
is therefore an aggregate property of the benchmark, not a statement about a typical
image, and a practitioner choosing a rule for a specific deployment cannot read it off
the pooled table: the evaluator changes which images fusion helps.

\paragraph{What of the index is fragile.} The same eight variants of Supplementary
Table~S6 were rescored on the same thirty images at the reported $400\times512$ grid, at
half and at double that grid, and at the reported grid after JPEG compression at
quality 50. The ordering is identical in all four conditions, SPGFusion first,
PIAFusion second and mean fusion last, while the values move: SPGFusion reads $0.799$,
$0.764$, $0.827$ and $0.787$. The instability of Sec.~\ref{sec:ranking} therefore belongs
to the construction of the index, its noise calibration and its width schedule, and not
to the resolution at which it is computed or to the storage of the fused image: two
configurations differing in one calibration constant reverse the first and the last of
eight methods, whereas changing the grid by a factor of two leaves the order untouched
and moves the values by up to $0.11$.

\paragraph{Matched sharpness, and the reporting threshold.} Two further controls
separate the rule's semantics from its arithmetic. Comparing rules at the same
temperature is unfair when their weight maps differ in spread, so we interpolate the
recorded curves to matched weight deviation: at a deviation of $0.03$ the three rules
are within $0.003$ mAP50 of each other, $0.8269$ to $0.8301$ under the visible-trained
evaluator, and at $0.06$ the top-hat prior is ahead of the activity rule by $0.0075$,
$0.8254$ against $0.8179$. The activity rule alone reaches the large deviations, where
it reads $0.8209$ under the visible-trained evaluator and $0.8728$ under the
modality-agnostic one. So part of the temperature sensitivity of the activity rule is
the sharpness of its weighting rather than the absence of a targetness map, but the
part that remains at matched sharpness is real and the prior keeps it.

The deficits also do not depend on the confidence at which detections are reported. At
floors of $0.05$, $0.25$ and $0.50$ the fused inputs stay below the visible source
under both YOLO families, SPGFusion by $0.178$, $0.181$ and $0.155$ F1 under YOLO11n,
and under the modality-agnostic detector the reversal shrinks as the floor rises, from
$+0.173$ for mean fusion at $0.05$ to $+0.032$ at $0.25$ and $-0.016$ at $0.50$. The
reversal of Sec.~\ref{sec:ranking} is therefore a property of the ranking of detections
rather than of the high-confidence subset, with one exception: under RT-DETR-l at a
floor of $0.05$ the source itself reaches an F1 of only $0.187$ and the fused inputs
gain slightly on it, $+0.012$ to $+0.015$.

\paragraph{What the targetness map is worth.} The variant set of
Sec.~\ref{sec:method} contains each rule and a scrambled copy of it, and the pair
isolates the contribution of the spatial correspondence between the weighting and the
image. Two results. First, the correspondence is what the index reads: destroying it
costs the reported index between $0.030$ and $0.110$ VIF depending on the rule and the
benchmark, with paired intervals that exclude zero in all twelve comparisons and
Benjamini--Hochberg adjusted $q$ values below $0.001$ (Supplementary Table~S11). Second,
a visible-trained detector does not read it: against its own scrambled copy, the
top-hat prior differs by $-0.0024$ $[-0.0113,+0.0053]$ and the learned probe by
$-0.0061$ $[-0.0138,+0.0015]$ mAP50 under the primary visible-trained detector, and the
weight-shuffled comparisons are equally flat, $+0.0068$ and $-0.0001$, so no interval
separates a rule from its own scrambled copy. Under the modality-agnostic detector the
correspondence becomes visible but small: the rules beat their scrambled copies by
$+0.0109$ to $+0.0153$ mAP50 with intervals that exclude zero. A fidelity gain of
$0.03$ to $0.11$ VIF therefore corresponds to nothing at all for one evaluator and to
about a hundredth of a mAP50 point for the other, which is the article's thesis at the
level of the rule rather than of the evaluator: the quantity that the metrics reward
is not the quantity that the detector uses.

\paragraph{Degraded inputs.} The negative result is a statement about clean inputs, and
the last experiment asks what happens when the visible modality is degraded. Every
input, the two unfused sources included, is degraded independently and identically
before scoring, with additive Gaussian noise ($\sigma=0.05$), Gaussian blur ($1.5$\,px)
and JPEG compression at quality 30, and each fused input is compared with the visible
source degraded in the same way. The picture changes with the degradation. Under
JPEG-30 and under blur the visible source loses $0.066$ and $0.138$ mAP50 and the
deficits shrink with it: three and four of the fifteen contrasts remain resolved, mean
fusion and the probe become level with the source ($-0.008$ to $+0.008$), and SPGFusion
still loses $0.082$ and $0.146$. Under additive noise the source collapses from $0.881$
to $0.530$ mAP50 and all fifteen fused inputs are resolved above it, by $+0.237$ to
$+0.299$ (Table~\ref{tab:corrupt}, Fig.~\ref{fig:corrupt}). We attribute this to what
makes the inputs different rather than to the rule: a fused image is a spatially
smoothed combination of two sources, so it carries less of the fine texture that
additive noise destroys, and averaging two independent noise realisations of the same
scene removes some of the noise itself. The experiment therefore bounds the claim of
Sec.~\ref{sec:detect} rather than extending it: the fused input is a worse input than
the unfused visible image when that image is intact, and can be a better one when it is
not, which is also the setting the field proposes fusion for. A split of the
degradation between the two arms, degrading the visible source alone, is what would
separate the two mechanisms, and we did not run it.

\begin{table*}[!t]\centering\small

\caption{Every input degraded, with the visible-trained detector: mAP50 of the
degraded visible source, the contrast of four rules against it, and how many of the
fifteen contrasts are resolved below it, unresolved and resolved above it. Every input
is degraded independently and identically.}\label{tab:corrupt}
\setlength{\tabcolsep}{4pt}
\begin{tabular*}{\linewidth}{@{\extracolsep{\fill}}l SSSS c@{}}
\toprule
Condition & {VIS source} & {mean} & {probe} & {SPGFusion} & {below/ns/above}\\
\midrule
clean & 0.8811 & -0.054 & -0.056 & -0.154 & 15 / 0 / 0\\
JPEG-30 & 0.8150 & +0.003 & +0.006 & -0.082 & 3 / 12 / 0\\
blur, $1.5$\,px & 0.7433 & -0.008 & +0.008 & -0.146 & 4 / 11 / 0\\
noise, $\sigma=0.05$ & 0.5296 & +0.298 & +0.299 & +0.237 & 0 / 0 / 15\\
\bottomrule
\end{tabular*}

\end{table*}

\begin{figure*}[!t]\centering

\includegraphics[width=\linewidth]{fig_corrupt_v2.pdf}
\caption{The negative result under degraded inputs. (a) Paired differences for four
rules under the clean, JPEG-30 and blur conditions on a common scale. (b) The same
rules under additive noise on its own scale because its effect is ten times larger and
of the opposite sign. (c) Each square is one of fifteen inputs: green, resolved below
the degraded visible source; grey, unresolved; orange, resolved above. The degraded
source scores are reported in Table~\ref{tab:corrupt}. Compression and blur shrink the
deficit without reversing it, whereas additive noise collapses the visible source and
puts every fused input above it.}
\label{fig:corrupt}

\end{figure*}

\paragraph{Which part of the rule carries it, and what it costs.} The three-axis
comparison separates the four rules at once. The targetness prior is ahead of
mean fusion by $0.088$ VIF on TNO, $0.100$ on MSRS and $0.111$ on RoadScene, but its
downstream margins, $0.0033$ mAP50 and $0.0045$ mIoU, are below the resolution of those
benchmarks: the paired floor is $0.030$ mAP50 over the 300 test images and $0.010$ mIoU
over the 361. Neither targetness rule therefore beats mean fusion in a downstream task
at the resolution these benchmarks offer, and the ensemble sits between the two rules
everywhere. The learned probe is level with the prior in fidelity on two benchmarks and
behind it in detection by $0.0048$, so what the map buys is fidelity and, under the
modality-agnostic evaluator, stability of the weighting rather than a downstream gain.
The cost is small: the probe runs in $19.1$\,ms at $512\times384$ with $0.07$\,M
parameters, against $538.5$\,ms and $75.86$\,M for the fusion network of the largest
end-to-end method in the comparison; its inference also requires DINO and CLIP feature
extraction and peaks at $1,455$\,MiB above the baseline (Table~\ref{tab:cost}).

\end{sloppypar}

\section{Task-Aware Adaptation}
\label{sec:adapt}

\begin{sloppypar}
If a fusion rule is scored by a downstream detector, the obvious remedy is to train the
rule against that detector. We test five variants: adapting the distilled probe against
a single evaluator, against a pool of evaluators, against a pool with two further seeds,
and adapting the whole network against the pooled detector loss, either by cycling
through the pool one evaluator per step (round robin) or by back-propagating the
largest loss of the pool (worst case). All variants are scored on the same images as
the untuned rule, with the resampling of one benchmark's images shared across its
evaluators so that the interval belongs to the average effect.

No variant raises the average. The best case is probe-only adaptation over a pool of
evaluators, whose change in average quality is $+0.0008$ with an interval of
$[-0.0045,+0.0078]$, that is zero to within eight thousandths of a mAP50 point, and
which repeats across seeds. That variant does move the rule between evaluators: it
gains $+0.0389$ on the visible-trained YOLO26n and loses $0.0114$, $0.0224$ and
$0.0188$ on three held-out evaluators, each with an interval that excludes zero, so
what it produces is a different allocation of the same quality rather than a better
rule. Whole-network adaptation is negative on both evaluator sets, at $-0.0183$ and
$-0.0361$ over the four first-benchmark evaluators and at $-0.1031$ for the pooled
round robin over the seven, with only the $(+0.012)$ and $(+0.003)$ first-benchmark
averages of the pooled probe and the round robin positive, which is what makes the
pooled loss an off-domain effect. The mechanism is visible in the images: adapting only
the probe leaves the synthesis almost where it was (drift SSIM $0.92$ to $0.94$ on both
benchmarks), whereas adapting the whole network rewrites it, and the rewrite is far
larger off-domain, with the mean absolute difference rising from $0.063$ on the first
benchmark to $0.154$ and $0.260$ on the second. The clearest single number is the
contrast of the output: the untuned rule produces FLIR images with a pixel standard
deviation of $0.1568$ and the worst-case variant $0.0967$, having lost nearly two
fifths of the contrast that makes a thermal target visible. Task-aware training
therefore buys no net gain in this setting, and the parameter-light rule trained once
for reconstruction and never for the task remains the best cross-evaluator choice we
found. Supplementary Sec.~23 and Tables S26 and S27 give the per-evaluator results and
the image drift behind this section.
\end{sloppypar}

\section{Discussion}
\label{sec:discussion}

\begin{sloppypar}
Three findings carry the article, and Fig.~\ref{fig:decision} places them in one
view. The first is that a fusion ranking is a property of the
evaluator and index before it is a property of the fused inputs. The same inputs are
reordered by training modality, architecture and index construction, and the last
reversal is complete. The method ranked first by the reported index comes last under a
reference construction of the same metric family. These changes recur across seeds and
runs and are therefore not explainable as sampling noise.

The second is the stable aggregate result for a visible-trained reference. No fused
input beats the unfused visible image at mAP50 on the independent set, and the 42
evaluator-input deficits range from $0.046$ to $0.246$. The pooled audit preserves
the sign with deficits from $0.014$ to $0.152$. The sign is set by the rule and the evaluator's
training modality, while the magnitude varies with architecture. Scoring the visible
source in the detector's colour convention and using the library's metric implementation
change the magnitude by at most $0.035$ mAP50 and do not change the sign on the pooled
set.

The third is the uncertainty boundary. Source alignment matters. In the best-aligned
clean-split quartile, mean fusion, the top-hat prior and the probe are above the
visible source in all four runs; in the worst-aligned quartile their estimates are
below it, and the published networks lose by $0.14$ to $0.17$. The alignment quality
co-varies with contrast and target size, so these subgroups do not isolate registration
causally. The pooled stratification is likewise evaluator-dependent. In the pooled
audit the deficit resolves for all fourteen rules in the small-target and
high-contrast strata, and for five to fourteen in the crowded stratum. That partition
does not transfer uniformly to the independent-test evaluators: the clean BH counts
are $1/1/0/3/2/0$ for YOLO11n, $13/14/11/14/14/14$ for YOLO26n and
$2/14/2/14/14/5$ for RT-DETR-l (Supplementary Table~S29). The strata are hypotheses
about where fusion may fail, not a universal mechanism.

\begin{figure*}[!t]\centering

\includegraphics[width=0.62\linewidth]{fig_decision_v2.pdf}
\caption{The measured ranking is conditional on the evaluation setting. The same fourteen
fused images and the same held-out split are fixed while two analytical choices vary:
the trained evaluator and the fidelity index. On the 300 held-out test images, all $42$
paired mAP50 contrasts under the three clean visible-trained evaluators lie below the
colour-scored visible source; the clean modality-agnostic YOLO11n reverses the sign for
four inputs in every run. Across the fourteen fused inputs, SPGFusion ranks first under
the reported VIF and last under the reference pixel-VIF construction. Counts use paired
mAP50 contrasts on the fixed 1,350/150/300 split; the test images were not used for
training, validation or checkpoint selection. Per-input intervals are in
Table~\ref{tab:cleansplit}, q values are in Supplementary Table~S28, and the index
endpoints are in
Table~\ref{tab:fidfull}. The three thumbnails show test frame 190002 and its
mean-fusion output only as an illustration.}
\label{fig:decision}

\end{figure*}

What follows for practice is therefore a choice, not only a reporting rule. A detector
trained on visible images has no pooled gain from fusion, so the default should be the
visible source. A deployment-specific check is still required when sources are unusually
well aligned or the target distribution differs: our best-aligned clean-split quartile
favours several fused rules. Fusion is justified when the detector is trained on thermal
or fused data: under the clean-split modality-agnostic detector four of the fourteen fused
inputs beat the visible source, and under the pooled detector ten do. When the evaluator
is not fixed at design time, it should be chosen first from the modality the deployment
actually sees, and the fusion rule then chosen for that evaluator and stratum rather than
on a pooled fidelity table.

Several limits bound these statements. The independent split now covers all three
visible-trained architectures and the modality-agnostic control, but only YOLO11n has
four clean runs; YOLO26n and RT-DETR-l have one clean run each. The pooled audit retains
single-channel fused inputs and is therefore reported only as a sensitivity analysis.
FLIR cannot be colour-rescored because its checkpoints are unavailable. The evidence is from
single-class pedestrian benchmarks, the adaptation study covers one backbone and one
pooling rule, and RT-DETR-l changes architecture and optimisation budget together.
Annotations are shared across modalities, so the experiment cannot exclude a labelling
convention that favours visible appearance. An independent ECC estimate found a median
source-pair residual of $2.53$ px and a best-quartile cutoff of $1.61$ px; the subgroup
analysis shows that residual registration and scene composition contribute materially to
the pooled deficit. The estimator is translational and descriptive, not ground-truth
alignment metadata, so it cannot isolate registration from the correlated changes in
contrast and target size.
\end{sloppypar}

\section{Conclusion}
\label{sec:conclusion}

\begin{sloppypar}
We evaluated whether image-fidelity metrics reliably rank visible-infrared fusion rules
for downstream pedestrian detection. The analysis holds the fused inputs fixed and
separates sampling, training, evaluator, metric-construction, source-convention and
alignment effects. It covers fourteen fused inputs, 26 detection checkpoints, one
segmentation checkpoint and two pedestrian-detection benchmarks.

The results show that the fidelity ranking is not a stable property of the fused images.
On the independent LLVIP split, all 42 contrasts under three visible-trained evaluators
are negative, with deficits from $0.046$ to $0.246$ mAP50. A modality-agnostic reference
reverses the sign for four inputs in every run, and the reported and reference VIF
constructions place SPGFusion first and last. Task-aware adaptation does not remove the
dependence. Whole-network adaptation loses nearly two fifths of the output contrast
off-domain, and no adapted variant raises the average beyond its uncertainty interval.

The practical implication is conditional. For a visible-trained detector, the visible
source should be the default input, and fidelity alone is insufficient evidence for
choosing a fusion rule. A fused input should be selected only after validation with the
deployment modality and on the relevant alignment, contrast and target-size strata.
These conclusions are bounded to pedestrian-detection benchmarks and the tested fusion
rules. The released per-image scores, fused images and analysis scripts make every
reported interval and ranking comparison recomputable.
\end{sloppypar}

\section*{Statements and Declarations}
\subsection*{Author contributions}
\textbf{Jianyong Zhao:} Conceptualization, Methodology, Software, Formal analysis,
Writing -- original draft. \textbf{Zigao Wu:} Supervision, Writing -- review \&
editing. \textbf{Mengfei Ma:} Investigation, Validation, Writing -- review \&
editing.

\subsection*{Competing interests}
The authors declare that they have no known competing financial interests or
personal relationships that could have appeared to influence the work reported in
this article.

\subsection*{Funding}
This research did not receive any specific grant from funding agencies in the
public, commercial, or not-for-profit sectors.

\subsection*{Data availability}
\begin{sloppypar}
The datasets used in this article are publicly available. TNO is distributed through
Figshare at \url{https://doi.org/10.6084/m9.figshare.1008029}, MSRS through
\url{https://github.com/Linfeng-Tang/Image-Fusion}, RoadScene through
\url{https://github.com/hanna-xu/RoadScene}, LLVIP through
\url{https://github.com/bupt-ai-cz/LLVIP}, and the FLIR ADAS thermal dataset that the
aligned second detection benchmark is derived from is distributed by Teledyne FLIR at
\url{https://www.flir.com/oem/adas/adas-dataset-form/}. The aligned release used here
pairs that dataset with the COCO annotations published at
\url{https://doi.org/10.5281/zenodo.13907794}, and the 1500 training and 300 test pairs
are the person-containing subset selected by the released script
\texttt{analysis/j7\_flir.py}. Each dataset is used under the licence stated in its own
distribution. The
download- and preparation scripts, together with the trained checkpoints, are provided
in the public repository.
\end{sloppypar}

\begin{sloppypar}
The analysis itself is released as a versioned deposit: the per-image detection score
files of every evaluator, the fused images of every rule at native resolution, the
aggregate tables behind every figure, and the scripts that rebuild them from the frozen
runs. It also contains the fixed 1,350/150/300 split manifest, the two one-run clean
training manifests and their checkpoint hashes. The deposit carries a SHA-256 manifest, and the archived release at
\url{https://doi.org/10.5281/zenodo.22984389} holds the scripts and the aggregate tables,
so a reader can recompute any interval in this article without retraining a detector.
\end{sloppypar}

\subsection*{Code availability}
The code for reproducing all experiments is available at
\url{https://github.com/lessThanbetter/ssgf} and is archived at
\url{https://doi.org/10.5281/zenodo.22984389} under an MIT licence. The repository
includes the models,
the trained checkpoints, the dataset setup script, the fusion code, the detection
and segmentation evaluation code, the scripts that produce the detection tables and
the paired intervals of Secs.~\ref{sec:detect} and~\ref{sec:adapt} (under
\texttt{analysis/}), and the fidelity-index comparisons of Supplementary Sec.~8, including
the rescoring script for the reference construction. The versioned analysis deposit
contains the per-image detection statistics used by the tables and intervals, the
fused-image outputs at native resolution, the aggregate results and a SHA-256 manifest.
The companion repository contains the analysis code and names the exact inputs each
script expects.

\subsection*{Declaration of generative AI and AI-assisted technologies in the writing process}
During the preparation of this article the authors used an AI-assisted language-editing
tool to improve the readability and language of the manuscript. After using this tool,
the authors reviewed and edited the content as needed and take full responsibility for
the content of the publication.

\begingroup
\footnotesize
\begin{thebibliography}{10}
\expandafter\ifx\csname url\endcsname\relax
  \def\url#1{\texttt{#1}}\fi
\expandafter\ifx\csname urlprefix\endcsname\relax\def\urlprefix{URL }\fi
\expandafter\ifx\csname href\endcsname\relax
  \def\href#1#2{#2} \def\path#1{#1}\fi

\bibitem{ma2019}
J.~Ma, Y.~Ma, C.~Li, {Infrared and visible image fusion methods and
  applications: A survey}, Information Fusion 45 (2019) 153--178.
\newblock \href {https://doi.org/10.1016/j.inffus.2018.02.004}
  {\path{doi:10.1016/j.inffus.2018.02.004}}.

\bibitem{gajjar2026visual}
T.~Gajjar and M.~Goyani, {Infrared--visible image fusion for robust visual
  perception: Methods, benchmarks, evaluation and open challenges}, The Visual
  Computer 42(11) (2026).
\newblock \href {https://doi.org/10.1007/s00371-026-04695-5}
  {\path{doi:10.1007/s00371-026-04695-5}}.

\bibitem{zhu2017}
P.~Zhu, X.~Ma, Z.~Huang, {Fusion of infrared-visible images using improved
  multi-scale top-hat transform and suitable fusion rules}, Infrared Physics \&
  Technology 81 (2017) 282--295.
\newblock \href {https://doi.org/10.1016/j.infrared.2017.01.013}
  {\path{doi:10.1016/j.infrared.2017.01.013}}.

\bibitem{zhang2021}
S.~Zhang, X.~Li, X.~Zhang, S.~Zhang, {Infrared and visible image fusion based
  on saliency detection and two-scale transform decomposition}, Infrared
  Physics \& Technology 114 (2021) 103626.
\newblock \href {https://doi.org/10.1016/j.infrared.2020.103626}
  {\path{doi:10.1016/j.infrared.2020.103626}}.

\bibitem{yin2013}
H.~Yin, S.~Li, L.~Fang, {Simultaneous image fusion and super-resolution using
  sparse representation}, Information Fusion 14~(3) (2013) 229--240.
\newblock \href {https://doi.org/10.1016/j.inffus.2012.01.008}
  {\path{doi:10.1016/j.inffus.2012.01.008}}.

\bibitem{nejati2015}
M.~Nejati, S.~Samavi, S.~Shirani, {Multi-focus image fusion using
  dictionary-based sparse representation}, Information Fusion 25 (2015) 72--84.
\newblock \href {https://doi.org/10.1016/j.inffus.2014.10.004}
  {\path{doi:10.1016/j.inffus.2014.10.004}}.

\bibitem{yin2015}
H.~Yin, {Sparse representation with learned multiscale dictionary for image
  fusion}, Neurocomputing 148 (2015) 600--610.
\newblock \href {https://doi.org/10.1016/j.neucom.2014.07.003}
  {\path{doi:10.1016/j.neucom.2014.07.003}}.

\bibitem{wang2014}
L.~Wang, B.~Li, L.-f. Tian, {Multi-modal medical image fusion using the
  inter-scale and intra-scale dependencies between image shift-invariant
  shearlet coefficients}, Information Fusion 19 (2014) 20--28.
\newblock \href {https://doi.org/10.1016/j.inffus.2012.03.002}
  {\path{doi:10.1016/j.inffus.2012.03.002}}.

\bibitem{zhang2016}
Q.~Zhang, X.~Maldague, {An adaptive fusion approach for infrared and visible
  images based on NSCT and compressed sensing}, Infrared Physics \& Technology
  74 (2016) 11--20.
\newblock \href {https://doi.org/10.1016/j.infrared.2015.11.003}
  {\path{doi:10.1016/j.infrared.2015.11.003}}.

\bibitem{wei2017}
C.~Wei, B.~Zhou, W.~Guo, {A three scale image transformation for infrared and
  visible image fusion}, 2017 20th International Conference on Information
  Fusion (Fusion) (2017) 1--6\newblock \href {https://doi.org/10.23919/icif.2017.8009748}
  {\path{doi:10.23919/icif.2017.8009748}}.

\bibitem{fu2009}
M.-Y. FU, C.~ZHAO, {FUSION OF INFRARED AND VISIBLE IMAGES BASED ON THE SECOND
  GENERATION CURVELET TRANSFORM}, JOURNAL OF INFRARED AND MILLIMETER WAVES
  28~(4) (2009) 254--258.
\newblock \href {https://doi.org/10.3724/sp.j.1010.2009.00254}
  {\path{doi:10.3724/sp.j.1010.2009.00254}}.

\bibitem{ma2021}
X.~Ma, Z.~Wang, S.~Hu, {Multi-focus image fusion based on multi-scale sparse
  representation}, Journal of Visual Communication and Image Representation 81
  (2021) 103328.
\newblock \href {https://doi.org/10.1016/j.jvcir.2021.103328}
  {\path{doi:10.1016/j.jvcir.2021.103328}}.

\bibitem{chen2007}
H.~Chen, P.~K. Varshney, {A human perception inspired quality metric for image
  fusion based on regional information}, Information Fusion 8~(2) (2007)
  193--207.
\newblock \href {https://doi.org/10.1016/j.inffus.2005.10.001}
  {\path{doi:10.1016/j.inffus.2005.10.001}}.

\bibitem{piella}
G.~Piella, H.~Heijmans, {A new quality metric for image fusion}, in: Proceedings
  of the 2003 International Conference on Image Processing (ICIP), vol.~3, 2003,
  pp.~173--176.
\newblock \href {https://doi.org/10.1109/ICIP.2003.1247209}
  {\path{doi:10.1109/ICIP.2003.1247209}}.

\bibitem{ma20191}
J.~Ma, W.~Yu, P.~Liang, C.~Li, J.~Jiang, {FusionGAN: A generative adversarial
  network for infrared and visible image fusion}, Information Fusion 48 (2019)
  11--26.
\newblock \href {https://doi.org/10.1016/j.inffus.2018.09.004}
  {\path{doi:10.1016/j.inffus.2018.09.004}}.

\bibitem{hou2021}
J.~Hou, D.~Zhang, W.~Wu, J.~Ma, H.~Zhou, {A Generative Adversarial Network for
  Infrared and Visible Image Fusion Based on Semantic Segmentation}, Entropy
  23~(3) (2021) 376.
\newblock \href {https://doi.org/10.3390/e23030376}
  {\path{doi:10.3390/e23030376}}.

\bibitem{yi2021}
S.~Yi, J.~Li, X.~Yuan, {DFPGAN: Dual fusion path generative adversarial network
  for infrared and visible image fusion}, Infrared Physics \& Technology 119
  (2021) 103947.
\newblock \href {https://doi.org/10.1016/j.infrared.2021.103947}
  {\path{doi:10.1016/j.infrared.2021.103947}}.

\bibitem{song2022}
A.~Song, H.~Duan, H.~Pei, L.~Ding, {Triple-discriminator generative adversarial
  network for infrared and visible image fusion}, Neurocomputing 483 (2022)
  183--194.
\newblock \href {https://doi.org/10.1016/j.neucom.2022.02.025}
  {\path{doi:10.1016/j.neucom.2022.02.025}}.

\bibitem{xu2022}
H.~Xu, J.~Ma, J.~Jiang, X.~Guo, H.~Ling, {U2Fusion: A Unified Unsupervised
  Image Fusion Network}, IEEE Transactions on Pattern Analysis and Machine
  Intelligence 44~(1) (2022) 502--518.
\newblock \href {https://doi.org/10.1109/tpami.2020.3012548}
  {\path{doi:10.1109/tpami.2020.3012548}}.

\bibitem{li2021}
H.~Li, X.-J. Wu, J.~Kittler, {RFN-Nest: An end-to-end residual fusion network
  for infrared and visible images}, Information Fusion 73 (2021) 72--86.
\newblock \href {https://doi.org/10.1016/j.inffus.2021.02.023}
  {\path{doi:10.1016/j.inffus.2021.02.023}}.

\bibitem{li2019}
H.~Li, X.-J. Wu, {DenseFuse: A Fusion Approach to Infrared and Visible Images},
  IEEE Transactions on Image Processing 28~(5) (2019) 2614--2623.
\newblock \href {https://doi.org/10.1109/tip.2018.2887342}
  {\path{doi:10.1109/tip.2018.2887342}}.

\bibitem{zhang2020}
Y.~Zhang, Y.~Liu, P.~Sun, H.~Yan, X.~Zhao, L.~Zhang, {IFCNN: A general image
  fusion framework based on convolutional neural network}, Information Fusion
  54 (2020) 99--118.
\newblock \href {https://doi.org/10.1016/j.inffus.2019.07.011}
  {\path{doi:10.1016/j.inffus.2019.07.011}}.

\bibitem{wang2024}
X.~Wang, L.~Fang, J.~Zhao, Z.~Pan, H.~Li, Y.~Li, {UUD-Fusion: An unsupervised
  universal image fusion approach via generative diffusion model}, Computer
  Vision and Image Understanding 249 (2024) 104218.
\newblock \href {https://doi.org/10.1016/j.cviu.2024.104218}
  {\path{doi:10.1016/j.cviu.2024.104218}}.

\bibitem{liu2023}
R.~Liu, H.~Wang, S.~Du, Y.~Liu, {A Visible and Infrared Image Fusion Framework
  Based on Dual-Path Encoder-Decoder and Multi-Scale Discrete Wavelet
  Transform}, 2023 IEEE International Conference on Image Processing (ICIP)
  (2023) 1995--1999\newblock \href {https://doi.org/10.1109/icip49359.2023.10222863}
  {\path{doi:10.1109/icip49359.2023.10222863}}.

\bibitem{hao2022}
S.~Hao, T.~He, B.~An, X.~Ma, H.~Wen, F.~Wang, {VDFEFuse: A novel fusion
  approach to infrared and visible images}, Infrared Physics \& Technology 121
  (2022) 104048.
\newblock \href {https://doi.org/10.1016/j.infrared.2022.104048}
  {\path{doi:10.1016/j.infrared.2022.104048}}.

\bibitem{yang2023}
X.~Yang, H.~Huo, R.~Wang, C.~Li, X.~Liu, J.~Li, {DGLT-Fusion: A decoupled
  global–local infrared and visible image fusion transformer}, Infrared
  Physics \& Technology 128 (2023) 104522.
\newblock \href {https://doi.org/10.1016/j.infrared.2022.104522}
  {\path{doi:10.1016/j.infrared.2022.104522}}.

\bibitem{an2024}
R.~An, G.~Liu, Y.~Qian, M.~Xing, H.~Tang, {MRASFusion: A multi-scale residual
  attention infrared and visible image fusion network based on semantic
  segmentation guidance}, Infrared Physics \& Technology 139 (2024) 105343.
\newblock \href {https://doi.org/10.1016/j.infrared.2024.105343}
  {\path{doi:10.1016/j.infrared.2024.105343}}.

\bibitem{yi2024}
X.~Yi, L.~Tang, H.~Zhang, H.~Xu, J.~Ma, {Diff-IF: Multi-modality image fusion
  via diffusion model with fusion knowledge prior}, Information Fusion 110
  (2024) 102450.
\newblock \href {https://doi.org/10.1016/j.inffus.2024.102450}
  {\path{doi:10.1016/j.inffus.2024.102450}}.

\bibitem{yang2025}
B.~Yang, Z.~Jiang, D.~Pan, H.~Yu, G.~Gui, W.~Gui, {LFDT-Fusion: A latent
  feature-guided diffusion Transformer model for general image fusion},
  Information Fusion 113 (2025) 102639.
\newblock \href {https://doi.org/10.1016/j.inffus.2024.102639}
  {\path{doi:10.1016/j.inffus.2024.102639}}.

\bibitem{tang2022}
L.~Tang, J.~Yuan, H.~Zhang, X.~Jiang, J.~Ma, {PIAFusion: A progressive infrared
  and visible image fusion network based on illumination aware}, Information
  Fusion 83-84 (2022) 79--92.
\newblock \href {https://doi.org/10.1016/j.inffus.2022.03.007}
  {\path{doi:10.1016/j.inffus.2022.03.007}}.

\bibitem{zhang2026}
H.~Zhang, S.~Yao, J.~Ma, J.~Jiang, Y.~Zhang, and H.~Zhou, {SPGFusion: Semantic
  prior guided infrared and visible image fusion via pretrained vision models},
  Information Fusion 125 (2026) 103433.
\newblock \href {https://doi.org/10.1016/j.inffus.2025.103433}
  {\path{doi:10.1016/j.inffus.2025.103433}}.

\bibitem{zhou2026}
Y.~Zhou, Q.~Yang, K.~Ren, Q.~Chen, {MTFusion: A semantic segmentation-guided
  infrared and visible image fusion network based on an attention-enhanced
  hybrid Mamba-Transformer architecture}, Infrared Physics \& Technology 153
  (2026) 106325.
\newblock \href {https://doi.org/10.1016/j.infrared.2025.106325}
  {\path{doi:10.1016/j.infrared.2025.106325}}.

\bibitem{bianco2020}
S.~Bianco, M.~Buzzelli, G.~Ciocca, R.~Schettini, {Neural architecture search
  for image saliency fusion}, Information Fusion 57 (2020) 89--101.
\newblock \href {https://doi.org/10.1016/j.inffus.2019.12.007}
  {\path{doi:10.1016/j.inffus.2019.12.007}}.

\bibitem{khan2025}
H.~Khan, M.~T. Usman, J.~Koo, {Bilateral Feature Fusion with hexagonal
  attention for robust saliency detection under uncertain environments},
  Information Fusion 121 (2025) 103165.
\newblock \href {https://doi.org/10.1016/j.inffus.2025.103165}
  {\path{doi:10.1016/j.inffus.2025.103165}}.

\bibitem{wang2025}
Q.~Wang, Z.~Li, S.~Zhang, Y.~Luo, W.~Chen, T.~Wang, N.~Chi, Q.~Dai,
  {SMAE-Fusion: Integrating saliency-aware masked autoencoder with hybrid
  attention transformer for infrared–visible image fusion}, Information
  Fusion 117 (2025) 102841.
\newblock \href {https://doi.org/10.1016/j.inffus.2024.102841}
  {\path{doi:10.1016/j.inffus.2024.102841}}.

\bibitem{xu2021}
H.~Xu, H.~Zhang, J.~Ma, {Classification Saliency-Based Rule for Visible and
  Infrared Image Fusion}, IEEE Transactions on Computational Imaging 7 (2021)
  824--836.
\newblock \href {https://doi.org/10.1109/tci.2021.3100986}
  {\path{doi:10.1109/tci.2021.3100986}}.

\bibitem{sheikh2006}
H.~R. Sheikh, A.~C. Bovik, {Image information and visual quality}, IEEE
  Transactions on Image Processing 15~(2) (2006) 430--444.
\newblock \href {https://doi.org/10.1109/TIP.2005.859378}
  {\path{doi:10.1109/TIP.2005.859378}}.

\bibitem{qin2020}
X.~Qin, Z.~Zhang, C.~Huang, M.~Dehghan, O.~R. Zaiane, M.~Jagersand, {U2-Net:
  Going deeper with nested U-structure for salient object detection}, Pattern
  Recognition 106 (2020) 107404.
\newblock \href {https://doi.org/10.1016/j.patcog.2020.107404}
  {\path{doi:10.1016/j.patcog.2020.107404}}.

\bibitem{zhu2025}
J.~Zhu, X.~Qin, A.~Elsaddik, {DC-Net: Divide-and-conquer for salient object
  detection}, Pattern Recognition 157 (2025) 110903.
\newblock \href {https://doi.org/10.1016/j.patcog.2024.110903}
  {\path{doi:10.1016/j.patcog.2024.110903}}.

\bibitem{kirillov2023sam}
A.~Kirillov, E.~Mintun, N.~Ravi, J.~Shlens, J.~Johnson, {Segment Anything}, in:
  Proceedings of the IEEE/CVF International Conference on Computer Vision
  (ICCV), 2023, pp. 4015--4026.
\newblock \href {https://doi.org/10.1109/ICCV51070.2023.00371}
  {\path{doi:10.1109/ICCV51070.2023.00371}}.

\bibitem{guo2024}
J.~Guo, Z.~Huang, Y.~Tao, {Multispectral Pedestrian Detection Based on
  Prior-Saliency Attention and Image Fusion}, Electronics 13~(9) (2024) 1770.
\newblock \href {https://doi.org/10.3390/electronics13091770}
  {\path{doi:10.3390/electronics13091770}}.

\bibitem{petilli2022}
E.~Petilli, {Monocular thermal imaging for pedestrian detection and ranging},
  Optical Engineering 62~(03) (2022) 031212.
\newblock \href {https://doi.org/10.1117/1.OE.62.3.031212}
  {\path{doi:10.1117/1.OE.62.3.031212}}.

\bibitem{li2010}
J.~Li, W.~Gong, W.~Li, X.~Liu, {Robust pedestrian detection in thermal infrared
  imagery using the wavelet transform}, Infrared Physics \& Technology 53~(4)
  (2010) 267--273.
\newblock \href {https://doi.org/10.1016/j.infrared.2010.03.005}
  {\path{doi:10.1016/j.infrared.2010.03.005}}.

\bibitem{redmon2016yolo}
J.~Redmon, S.~Divvala, R.~Girshick, A.~Farhadi, {You Only Look Once: Unified,
  Real-Time Object Detection}, in: Proceedings of the IEEE Conference on
  Computer Vision and Pattern Recognition (CVPR), 2016, pp. 779--788.
\newblock \href {https://doi.org/10.1109/CVPR.2016.91}
  {\path{doi:10.1109/CVPR.2016.91}}.

\bibitem{jocher2024yolo}
G.~Jocher, J.~Qiu, {Ultralytics YOLO}, software release (2024).
\newblock \href {https://github.com/ultralytics/ultralytics}
  {\path{https://github.com/ultralytics/ultralytics}}.

\bibitem{lv2024rtdetr}
W.~Lv, S.~Xu, Y.~Zhao, G.~Wang, J.~Wei, C.~Cui, Y.~Du, Q.~Dang, Y.~Liu,
  {DETRs Beat YOLOs on Real-time Object Detection}, in: Proceedings of the
  IEEE/CVF Conference on Computer Vision and Pattern Recognition (CVPR), 2024,
  pp. 16965--16974.
\newblock \href {https://doi.org/10.1109/CVPR52733.2024.01605}
  {\path{doi:10.1109/CVPR52733.2024.01605}}.

\bibitem{ronneberger2015}
O.~Ronneberger, P.~Fischer, T.~Brox, {U-Net: Convolutional Networks for
  Biomedical Image Segmentation}, Lecture Notes in Computer Science 9351 (2015)
  234--241\newblock \href {https://doi.org/10.1007/978-3-319-24574-4_28}
  {\path{doi:10.1007/978-3-319-24574-4_28}}.

\bibitem{flir_adas}
Teledyne FLIR, {FLIR Thermal Dataset for Algorithm Training}, dataset (2018).
\newblock \href {https://www.flir.com/oem/adas/adas-dataset-form/}
  {\path{https://www.flir.com/oem/adas/adas-dataset-form/}}.

\end{thebibliography}

\endgroup

\end{document}
